% ============================================================
%  Géolocalisation des Eglises EEC — ENSPY 2025-2026
%  pdfLaTeX + biber
%  Compilation : pdflatex → biber → pdflatex → pdflatex
% ============================================================
\documentclass[12pt, a4paper]{report}

\usepackage[utf8]{inputenc}
\usepackage[T1]{fontenc}
\usepackage[french]{babel}
\usepackage{csquotes}
\usepackage[a4paper, top=2.5cm, bottom=2.5cm,
            left=3cm, right=2.5cm, headheight=14pt]{geometry}
\usepackage{lmodern}
\usepackage{microtype}
\usepackage{setspace}
\setlength{\parindent}{1.2cm}
\setlength{\parskip}{0pt}

% ---- Couleurs ----
\usepackage{xcolor}
\definecolor{eecgreen}{RGB}{27,94,32}
\definecolor{eccdarkgreen}{RGB}{17,60,20}
\definecolor{eecyellow}{RGB}{249,168,37}
\definecolor{eeclightgreen}{RGB}{232,245,233}
\definecolor{bordergold}{RGB}{212,160,23}

% ---- Bordure double sur chaque page ----
\usepackage{eso-pic}
\usepackage{tikz}
\AddToShipoutPictureBG{%
  \begin{tikzpicture}[remember picture, overlay]
    \draw[color=bordergold, line width=2pt]
      ([xshift=10pt,yshift=-10pt]current page.north west)
      rectangle
      ([xshift=-10pt,yshift=10pt]current page.south east);
    \draw[color=eecgreen, line width=0.7pt]
      ([xshift=15pt,yshift=-15pt]current page.north west)
      rectangle
      ([xshift=-15pt,yshift=15pt]current page.south east);
  \end{tikzpicture}}

% ---- En-têtes et pieds — texte noir ----
\usepackage{fancyhdr}
\pagestyle{fancy}
\fancyhf{}
\fancyhead[L]{\small\textit{Géolocalisation des Eglises EEC --- ENSPY 2025-2026}}
\fancyhead[R]{\small\thepage}
\fancyfoot[C]{\small\textit{Département de Génie Informatique}}
\renewcommand{\headrulewidth}{0.4pt}
\renewcommand{\footrulewidth}{0pt}

% ---- Titres de chapitre ----
\usepackage{titlesec}
% Chapitres numérotés — boîte décorative verte + numéro + titre en petites capitales
\titleformat{\chapter}[display]{%
  \normalfont%
}{%
  \begin{tikzpicture}[remember picture,overlay]%
    \node[rotate=90, anchor=center,
          text=eecgreen!80!black,
          font=\footnotesize\bfseries\scshape,
          inner sep=0pt]
      at ([xshift=-72pt, yshift=-55pt]current page.north east)
      {Chapitre};%
    \fill[eecgreen]
      ([xshift=-60pt, yshift=-10pt]current page.north east)
      rectangle
      ([xshift=0pt, yshift=-100pt]current page.north east);%
    \node[text=white, font=\fontsize{68}{68}\selectfont\bfseries, anchor=center]
      at ([xshift=-30pt, yshift=-55pt]current page.north east)
      {\thechapter};%
  \end{tikzpicture}%
}{10pt}{%
  \centering\normalfont\large\bfseries\scshape\color{eecgreen}%
}[\vspace{2pt}{\color{eecyellow}\hrule height 2.5pt}\vspace{4pt}]
\titlespacing{\chapter}{0pt}{26pt}{14pt}

% Chapitres non numérotés (\chapter*) — style simple sans décoration
\titleformat{name=\chapter,numberless}[block]{%
  \normalfont\Large\bfseries\color{eecgreen}%
}{}{0pt}{\MakeUppercase}[%
  \vspace{3pt}{\color{eecyellow}\hrule height 2.5pt}\vspace{5pt}%
]
\titlespacing{name=\chapter,numberless}{0pt}{0pt}{10pt}

% Sections et sous-sections — hiérarchie verte grasse
\titleformat{\section}
  {\normalfont\large\bfseries\color{eccdarkgreen}}
  {\thesection.}{8pt}{}
  [\vspace{1pt}{\color{eecgreen!40!white}\hrule height 0.6pt}\vspace{2pt}]
\titleformat{\subsection}
  {\normalfont\normalsize\bfseries\color{eecgreen}}
  {\thesubsection.}{7pt}{}
\titleformat{\subsubsection}
  {\normalfont\normalsize\bfseries\color{eecgreen!90!black}}
  {\thesubsubsection.}{6pt}{}

% ---- Tableaux ----
\usepackage{array}
\usepackage{colortbl}
\usepackage{booktabs}
\usepackage{longtable}

% ---- Figures et flottants ----
\usepackage{graphicx}
\usepackage{float}

% ---- Listes ----
\usepackage{enumitem}

% ---- Lettrine (lettre ornée de début de chapitre) ----
\usepackage{lettrine}
\renewcommand{\LettrineFontHook}{\color{eecgreen}\bfseries}
\setcounter{DefaultLines}{2}
\setlength{\DefaultFindent}{4pt}
\setlength{\DefaultNfindent}{0pt}

% ---- Symboles et mathématiques ----
\usepackage{pifont}
\usepackage{amsmath}
\usepackage{amssymb}

% ---- Liens ----
\usepackage[colorlinks=true, linkcolor=eecgreen,
            urlcolor=eecgreen, citecolor=eecgreen]{hyperref}

% ---- Sommaire local par chapitre (minitoc — doit être APRÈS hyperref) ----
\usepackage{minitoc}
\mtcsettitle{minitoc}{\normalsize\bfseries\color{eecgreen}Sommaire}
\setcounter{minitocdepth}{3}

% ---- Symboles tableaux comparatifs ----
\newcommand{\oui}{\textcolor{eecgreen}{\ding{51}}}
\newcommand{\non}{\textcolor{red}{\ding{55}}}
\newcommand{\partiel}{\textcolor{orange}{$\sim$}}

% ---- Mots-clés techniques (premier emploi) ----
\newcommand{\kw}[1]{\textbf{#1}}

% ---- Plan de chapitre (boîte colorée en tête de chapitre) ----
\newenvironment{chapplan}{%
  \vspace{6pt}%
  \noindent\colorbox{eeclightgreen!70}{%
    \begin{minipage}[t]{\dimexpr\linewidth-2\fboxsep}%
      \vspace{5pt}%
      \noindent{\small\bfseries\color{eccdarkgreen}%
        \ding{228}\quad Plan du chapitre}%
      \vspace{3pt}\\%
      {\color{eecgreen!60!white}\hrule height 0.5pt}%
      \vspace{2pt}%
      \begin{itemize}[%
        leftmargin=1.4cm, itemsep=1pt, topsep=2pt, parsep=0pt,%
        label={\small\color{eecgreen}\ding{118}}]%
      \small%
}{%
      \end{itemize}%
      \vspace{4pt}%
    \end{minipage}%
  }%
  \vspace{8pt}%
}

% ---- Fiche de cas d'utilisation (bordure verte, en-tête coloré) ----
\newenvironment{ucfiche}[2]{%
  \vspace{6pt}%
  \noindent\colorbox{eccdarkgreen}{\parbox{\dimexpr\linewidth-2\fboxsep}{%
    \color{white}\bfseries\small #1~--- #2}}\par
  \noindent\fcolorbox{eecgreen!40!white}{white}{%
    \begin{minipage}[t]{\dimexpr\linewidth-2\fboxsep-2\fboxrule}%
      \vspace{4pt}%
      \small
}{%
      \vspace{2pt}%
    \end{minipage}%
  }%
  \vspace{6pt}%
}

% ============================================================
\begin{document}
\pagenumbering{roman}
\setstretch{1.3}

% ============================================================
%  RÉSUMÉ
% ============================================================
\chapter*{Résumé}
\addcontentsline{toc}{chapter}{Résumé}
\markboth{Résumé}{Résumé}

\noindent
L'Église Évangélique du Cameroun regroupe 858~paroisses, 311~œuvres sociales
et éducatives, ainsi que 685~ouvriers ecclésiastiques répartis dans 22~régions
synodales et 159~districts couvrant l'ensemble du territoire national, sans
que l'institution ne dispose d'aucun outil numérique centralisé permettant
leur visualisation géographique ni leur administration hiérarchisée. Cette
lacune compromet la capacité du Bureau National à piloter ses ressources
territoriales, à produire des statistiques consolidées et à assurer une mise
à jour contrôlée des données. Ce travail présente la conception et le
développement d'une plateforme web cartographique institutionnelle répondant
à ces besoins. L'approche adoptée s'appuie sur une architecture à services
séparés~: un serveur d'application gère les données et la logique métier, une
interface web assure la consultation et l'administration, une base de données
spécialisée dans les informations géographiques stocke les positions des
entités, un serveur cartographique diffuse les cartes selon des normes
internationales reconnues, et une bibliothèque de cartographie interactive
affiche les données côté navigateur~; l'ensemble est déployé sous une forme
conteneurisée, reproductible sur n'importe quel serveur. La plateforme ainsi
développée offre une carte interactive avec regroupement automatique des
marqueurs, des filtres hiérarchiques par région, district et type d'entité,
un panneau de détail par entité, ainsi qu'un module d'administration
multi-niveaux fondé sur un contrôle d'accès selon le rôle de chaque
utilisateur, à quatre profils~: national, régional, district et paroissial.
Par rapport aux solutions existantes analysées, ce système se distingue par
son alignement strict avec le référentiel hiérarchique synodal de l'Église,
la diffusion normalisée des données géographiques, et un déploiement
entièrement conteneurisé garantissant la reproductibilité de l'environnement
d'exécution.

\smallskip
\noindent\textbf{Mots-clés~:}
géolocalisation, système d'information géographique, Église Évangélique du
Cameroun, cartographie interactive, contrôle d'accès par rôle.

\vspace{0.8cm}
\noindent{\color{eecgreen}\rule{\linewidth}{0.8pt}}
\vspace{0.3cm}

% ============================================================
%  ABSTRACT
% ============================================================
\chapter*{Abstract}
\addcontentsline{toc}{chapter}{Abstract}
\markboth{Abstract}{Abstract}

\noindent
The Evangelical Church of Cameroon encompasses 858~parishes, 311~social and
educational institutions, and 685~ordained workers spread across
22~synodal regions and 159~districts covering the entire national territory,
yet the institution lacks any centralized digital system for their geographic
visualization or hierarchical administration. This gap undermines the National
Bureau's ability to monitor its territorial resources, generate consolidated
statistics, and maintain controlled data updates. This work presents the design
and development of an institutional cartographic web platform addressing these
needs. The approach relies on a separated service architecture~: an application
server handles the business logic and data, a web interface provides
consultation and administration, a dedicated geographic database stores the
location of each entity, a cartographic server publishes the maps according to
recognized international standards, and an interactive mapping library displays
the data in the browser~; the whole system is deployed in a containerized,
reproducible form that can run on any server. The resulting platform delivers
an interactive map with automatic marker clustering, hierarchical filters by
region, district, and entity type, a per-entity detail panel, and a
multi-level administration module enforcing access control based on each
user's role, across four profiles~: national, regional, district, and parish.
Compared to the existing solutions reviewed in this work, the system is
distinguished by its strict alignment with the Church's synodal hierarchical
framework, the standards-compliant publication of geographic data, and a
fully containerized deployment ensuring environment reproducibility.

\smallskip
\noindent\textbf{Keywords~:}
geolocation, geographic information system, Evangelical Church of Cameroon,
interactive cartography, role-based access control.

\clearpage

% ============================================================
%  SIGLES ET ABRÉVIATIONS
% ============================================================
\chapter*{Sigles et Abréviations}
\addcontentsline{toc}{chapter}{Sigles et Abréviations}
\markboth{Sigles et Abréviations}{Sigles et Abréviations}

\noindent
Le tableau ci-dessous recense par ordre alphabétique les sigles et abréviations
employés dans ce document.

\vspace{4pt}

\begingroup
\setstretch{1.0}
\renewcommand{\arraystretch}{1.12}
\noindent
\begin{tabular}{%
  >{\bfseries}p{3.2cm}%
  p{10.8cm}}
  \rowcolor{eecgreen}
  \multicolumn{1}{>{\bfseries\color{white}}p{3.2cm}}{Sigle} &
  \multicolumn{1}{>{\bfseries\color{white}}p{10.8cm}}{Signification complète} \\
  \rowcolor{eeclightgreen!60}
  API     & Application Programming Interface \\
  CORS    & Cross-Origin Resource Sharing \\
  \rowcolor{eeclightgreen!60}
  CRUD    & Create, Read, Update, Delete \\
  CSS     & Cascading Style Sheets \\
  \rowcolor{eeclightgreen!60}
  DRF     & Django REST Framework \\
  EEC     & Église Évangélique du Cameroun \\
  \rowcolor{eeclightgreen!60}
  ENSPY   & École Nationale Supérieure Polytechnique de Yaoundé \\
  GIS     & Geographic Information System \\
  \rowcolor{eeclightgreen!60}
  GPS     & Global Positioning System (Système de Positionnement Global) \\
  HTML    & HyperText Markup Language \\
  \rowcolor{eeclightgreen!60}
  HTTP    & HyperText Transfer Protocol \\
  HTTPS   & HyperText Transfer Protocol Secure \\
  \rowcolor{eeclightgreen!60}
  JSON    & JavaScript Object Notation \\
  JWT     & JSON Web Token \\
  \rowcolor{eeclightgreen!60}
  OGC     & Open Geospatial Consortium \\
  ORM     & Object-Relational Mapping \\
  \rowcolor{eeclightgreen!60}
  PDF     & Portable Document Format \\
  PostGIS & PostgreSQL Spatial Extension \\
  \rowcolor{eeclightgreen!60}
  RBAC    & Role-Based Access Control (Contrôle d'accès basé sur les rôles) \\
  REST    & Representational State Transfer \\
  \rowcolor{eeclightgreen!60}
  SIG     & Système d'Information Géographique \\
  SLD     & Styled Layer Descriptor \\
  \rowcolor{eeclightgreen!60}
  SND30   & Stratégie Nationale de Développement 2020--2030 \\
  SQL     & Structured Query Language \\
  \rowcolor{eeclightgreen!60}
  SRS     & Spatial Reference System \\
  SSL     & Secure Sockets Layer \\
  \rowcolor{eeclightgreen!60}
  SSR     & Server-Side Rendering \\
  TOTP    & Time-based One-Time Password \\
  \rowcolor{eeclightgreen!60}
  UI      & User Interface (Interface utilisateur) \\
  URL     & Uniform Resource Locator \\
  \rowcolor{eeclightgreen!60}
  UX      & User Experience (Expérience utilisateur) \\
  WFS     & Web Feature Service \\
  \rowcolor{eeclightgreen!60}
  WGS84   & World Geodetic System 1984 \\
  WMS     & Web Map Service \\
  \rowcolor{eeclightgreen!60}
  WSL     & Windows Subsystem for Linux \\
\end{tabular}
\endgroup

\clearpage

% ============================================================
%  GLOSSAIRE
% ============================================================
\chapter*{Glossaire}
\addcontentsline{toc}{chapter}{Glossaire}
\markboth{Glossaire}{Glossaire}

\noindent
Ce glossaire définit les termes techniques spécialisés employés dans ce document.

\vspace{6pt}

\noindent\textbf{Cluster de marqueurs~:}
Technique de regroupement visuel sur une carte qui agrège les marqueurs
géographiquement proches en un seul indicateur numérique. Lorsque l'utilisateur
zoome, le cluster éclate pour révéler les marqueurs individuels. Cette technique
améliore la lisibilité d'une carte comportant un grand nombre de points.

\medskip
\noindent\textbf{Conteneurisation~:}
Approche de déploiement logiciel qui isole une application et toutes ses
dépendances dans une unité standardisée appelée conteneur. Dans ce projet, Docker
et Docker Compose orchestrent l'ensemble des services (backend, base de données,
serveur cartographique, cache, frontend) de manière reproductible et indépendante
de l'environnement hôte.

\medskip
\noindent\textbf{GeoServer~:}
Serveur cartographique open source conforme aux standards de l'Open Geospatial
Consortium (OGC). Il permet de publier, gérer et diffuser des données
géographiques sous forme de services WMS et WFS, accessibles depuis n'importe
quel client cartographique compatible.

\medskip
\noindent\textbf{Géolocalisation~:}
Processus de détermination et d'enregistrement de la position géographique d'une
entité à l'aide de coordonnées GPS (latitude, longitude). Dans ce projet, la
géolocalisation désigne l'association de chaque paroisse et œuvre de l'EEC à des
coordonnées précises sur le territoire camerounais.

\medskip
\noindent\textbf{PostGIS~:}
Extension spatiale pour le système de gestion de base de données relationnelle
PostgreSQL. PostGIS ajoute la prise en charge des types de données géographiques
(points, lignes, polygones) et des fonctions de calcul spatial (distances,
intersections, projections). Il constitue le standard de facto pour les bases de
données spatiales open source.

\medskip
\noindent\textbf{RBAC --- Contrôle d'accès basé sur les rôles~:}
Modèle de sécurité informatique dans lequel les droits d'accès aux ressources
d'un système sont déterminés par le rôle attribué à chaque utilisateur, et non
par son identité propre. Quatre rôles sont définis~: national, régional, district
et paroissial, chacun disposant d'un périmètre d'action délimité.

\medskip
\noindent\textbf{Région synodale~:}
Entité administrative de l'Église Évangélique du Cameroun regroupant un ensemble
de districts et de paroisses dans une zone géographique définie. L'EEC compte
22~régions synodales couvrant l'ensemble du territoire camerounais. La région
synodale est le premier niveau hiérarchique sous le Bureau National.

\medskip
\noindent\textbf{Shapefile~:}
Format de fichier géospatial vectoriel développé par ESRI, largement utilisé pour
la représentation de données géographiques. Un shapefile est constitué de plusieurs
fichiers complémentaires (\texttt{.shp}, \texttt{.dbf}, \texttt{.prj},
\texttt{.shx}) qui encodent la géométrie, les attributs et le système de
coordonnées des entités géographiques.

\medskip
\noindent\textbf{SLD --- Styled Layer Descriptor~:}
Standard OGC permettant de décrire la symbolisation et l'apparence d'une couche
géographique publiée via WMS. Les fichiers SLD définissent les couleurs, tailles
et règles d'affichage des entités géographiques sur la carte.

\medskip
\noindent\textbf{Tuile cartographique~:}
Image raster de dimensions fixes (généralement 256~$\times$~256~pixels)
représentant une portion de carte à un niveau de zoom donné. Les bibliothèques
cartographiques comme Leaflet assemblent dynamiquement ces tuiles pour former la
carte visible par l'utilisateur.

\medskip
\noindent\textbf{WFS --- Web Feature Service~:}
Standard OGC permettant de requêter et de récupérer des données géographiques
vectorielles (points, lignes, polygones) depuis un serveur cartographique,
généralement au format GeoJSON ou GML. Contrairement au WMS qui retourne des
images, le WFS retourne des données brutes manipulables côté client.

\medskip
\noindent\textbf{WMS --- Web Map Service~:}
Standard OGC permettant de requêter des cartes sous forme d'images raster depuis
un serveur cartographique. Le client envoie une requête HTTP spécifiant l'étendue
géographique et les couches souhaitées~; le serveur retourne une image PNG ou
JPEG correspondante, prête à être superposée sur la carte.

\clearpage

% ============================================================
%  TABLE DES MATIÈRES
% ============================================================
\dominitoc
\tableofcontents
\clearpage

% ============================================================
%  LISTE DES FIGURES
% ============================================================
\listoffigures
\addcontentsline{toc}{chapter}{Liste des figures}
\clearpage

% ============================================================
%  LISTE DES TABLEAUX
% ============================================================
\listoftables
\addcontentsline{toc}{chapter}{Liste des tableaux}
\clearpage

% ============================================================
%  CORPS DU RAPPORT — numérotation arabe
% ============================================================
\pagenumbering{arabic}
\setcounter{page}{1}

% ============================================================
%  INTRODUCTION GÉNÉRALE
% ============================================================
\chapter*{Introduction Générale}
\addcontentsline{toc}{chapter}{Introduction Générale}
\markboth{Introduction Générale}{Introduction Générale}

La transformation numérique des institutions publiques et religieuses constitue
l'un des axes stratégiques du développement durable en Afrique subsaharienne.
Selon l'Union Internationale des Télécommunications, le taux de pénétration
d'Internet en Afrique subsaharienne a atteint 36~\% en 2023~[1],
progressant de plus de dix points en cinq ans grâce au déploiement massif des
réseaux mobiles. La région comptait en 2023 plus de 615~millions d'abonnés à la
téléphonie mobile~[2], une dynamique qui confère aux technologies web
un rôle central dans la modernisation des organisations de toute nature. Au
Cameroun, cette tendance est amplifiée par la Stratégie Nationale de
Développement 2020--2030 (SND30), qui identifie explicitement la numérisation de
l'administration publique et des institutions nationales comme l'un des leviers
prioritaires du développement économique du pays~[3].

Dans ce contexte, les Systèmes d'Information Géographique (SIG) se sont imposés
comme des outils incontournables pour la gestion et la visualisation des données
territoriales. Longley et al.~définissent un SIG comme un système informatique
conçu pour la collecte, le stockage, la manipulation, l'analyse et la
représentation de données géoréférencées~[4]. À l'échelle
mondiale, ces plateformes sont déployées dans des domaines aussi variés que la
santé publique, la gestion urbaine, la logistique humanitaire et la planification
administrative. La standardisation des échanges de données géographiques par
l'Open Geospatial Consortium (OGC) --- notamment les protocoles Web Map Service
(WMS)~[5] et Web Feature Service (WFS)~[6] --- a
rendu ces technologies accessibles à des organisations de toute taille et de tout
secteur, y compris les institutions confessionnelles. En Afrique centrale, les
organisations religieuses constituent souvent le premier réseau d'établissements
scolaires, de centres de santé et de structures d'encadrement rural présent dans
les zones éloignées des grandes agglomérations~[7]. Leur capacité
à gérer efficacement des ressources géographiquement dispersées à travers de
larges territoires nationaux constitue un enjeu direct pour le développement
des communautés qu'elles desservent.

L'Église Évangélique du Cameroun (EEC), fondée en 1957 dans la continuité des
missions protestantes bâloises et parisiennes, est l'une des plus grandes
institutions protestantes d'Afrique centrale. Son organisation territoriale
repose sur une hiérarchie administrative à quatre niveaux~: le Bureau National,
qui assure la direction centrale de l'institution~; les 22~Régions Synodales, qui
couvrent l'ensemble du territoire camerounais~; les 159~districts ecclésiaux, qui
coordonnent les activités locales~; et les 858~paroisses, qui constituent le
maillage de base de la présence de l'EEC sur le terrain. À ces entités s'ajoutent
311~œuvres sociales et éducatives --- établissements scolaires, centres médicaux,
exploitations agropastorales, immeubles et terrains --- et plus de
685~ouvriers ecclésiastiques (pasteurs, évangélistes, diacres permanents) répartis
sur l'ensemble de ce réseau.

Malgré cette envergure nationale considérable, l'EEC ne dispose à ce jour d'aucun
système numérique centralisé permettant la visualisation géographique de ses
entités, leur administration hiérarchisée selon les périmètres synodaux, ni la
consolidation automatique de leurs données statistiques. La localisation des
paroisses, la gestion des affectations des ouvriers et le suivi des statistiques
annuelles sont assurés manuellement via des fichiers tableur transmis
périodiquement entre niveaux hiérarchiques. Cette situation génère plusieurs
problèmes opérationnels concrets~: l'absence de vue géographique d'ensemble
empêche le Bureau National d'identifier les zones sous-dotées en ressources
pastorales ou d'optimiser l'affectation des ouvriers~; la dispersion des données
dans des fichiers non interconnectés rend toute consolidation statistique nationale
laborieuse et sujette aux erreurs~; l'absence de contrôle d'accès formalisé
compromet l'intégrité des informations institutionnelles~; et l'audit des données
sources révèle que 480~paroisses, soit 56~\% de l'effectif total, ne disposent
d'aucune coordonnée GPS valide, ce qui illustre l'inexistence d'un processus
structuré de collecte et de validation des données de localisation.

Plusieurs catégories de solutions pourraient théoriquement répondre à ces besoins.
Les plateformes SIG généralistes comme ArcGIS~Online~[8] ou
Google~Maps~Platform~[9] offrent des capacités cartographiques
avancées mais demeurent des outils génériques, soumis à des licences commerciales
onéreuses et dépourvus de toute logique hiérarchique institutionnelle. Les outils
open~source comme QGIS~[10] ou OpenLayers permettent la visualisation
cartographique avancée mais ne constituent pas des plateformes web multi-utilisateurs
dotées de contrôle d'accès et de gestion de données métier. Une première version
d'une plateforme de géolocalisation EEC avait été développée antérieurement --- elle
offrait une cartographie de base mais reposait sur des fichiers GeoJSON statiques,
ne disposait pas d'un mécanisme d'authentification par sessions sécurisées, et son
contrôle d'accès ne correspondait pas à la granularité de la hiérarchie synodale.
Aucune des solutions recensées ne répond simultanément aux exigences d'une
plateforme institutionnelle sécurisée, intégrée, alignée sur la structure synodale
de l'EEC et déployable dans un environnement à infrastructure informatique limitée.

\medskip

Face à ces constats, la problématique suivante guide ce travail~:

\vspace{6pt}
\noindent\colorbox{eeclightgreen!80}{\parbox{0.88\linewidth}{%
  \vspace{6pt}%
  \noindent\textit{\textbf{Question de recherche~:} Comment concevoir et
  implémenter une plateforme web cartographique institutionnelle répondant
  simultanément aux exigences de visualisation géographique interactive,
  d'administration hiérarchisée à quatre niveaux de l'autorité synodale et de
  sécurité des données de l'Église Évangélique du Cameroun~?}%
  \vspace{6pt}%
}}
\vspace{6pt}

Nous posons l'hypothèse qu'une architecture à services découplés --- combinant un
backend Django~5 avec GeoDjango et Django REST Framework pour la logique métier et
les interfaces de programmation (API), une base de données PostgreSQL avec
l'extension PostGIS pour la persistance des données spatiales, un serveur
cartographique GeoServer~2.26 pour la publication des couches géographiques selon
les standards OGC (WMS et WFS), et une interface Next.js~16 pour la consultation
publique et l'administration --- permettrait de répondre à cette problématique. La
sécurité reposerait sur un mécanisme d'authentification par sessions Django
(\textit{HttpOnly}, \textit{Secure}, \textit{SameSite}) et un contrôle d'accès
basé sur les rôles (RBAC) à quatre niveaux hiérarchiques aligné sur la structure
synodale. La conteneurisation de l'ensemble via Docker Compose garantirait la
reproductibilité de l'environnement d'exécution et la facilité du déploiement.

\medskip

Pour répondre à cette problématique, ce travail se fixe les objectifs suivants~:

\begin{enumerate}[leftmargin=1.5cm, itemsep=2pt, topsep=4pt]
  \item \textbf{Modéliser} la structure de données hiérarchique de l'EEC
    (régions synodales, districts, paroisses, œuvres sociales, ouvriers
    ecclésiastiques) dans un schéma de base de données relationnelle-spatiale
    PostGIS intégrant les contraintes d'intégrité référentielle et les
    informations de géolocalisation~;

  \item \textbf{Développer} un backend API REST sécurisé avec Django~5 et
    Django REST Framework, implémentant un contrôle d'accès basé sur les rôles
    à quatre niveaux --- national, régional, district et paroissial --- et un
    journal d'audit automatique de toutes les opérations d'écriture~;

  \item \textbf{Configurer} un serveur cartographique GeoServer~2.26 publiant
    les couches géographiques de l'EEC selon les standards OGC WMS~1.3.0 et
    WFS~2.0.0, avec des styles SLD aux couleurs institutionnelles de l'EEC~;

  \item \textbf{Implémenter} une interface cartographique interactive basée sur
    Leaflet.js avec regroupement automatique des marqueurs (\textit{clustering}),
    filtres hiérarchiques multicritères, popups de détail par entité et contrôle
    dynamique des couches cartographiques~;

  \item \textbf{Réaliser} un module d'administration multi-niveaux en Next.js~16
    permettant le CRUD des entités selon le périmètre du rôle, l'import des
    données depuis les fichiers sources, l'export PDF et Excel, et la
    visualisation des statistiques avancées par région et par année~;

  \item \textbf{Conteneuriser} l'ensemble de la plateforme avec Docker Compose
    en cinq services afin de garantir la reproductibilité de l'environnement
    d'exécution et de simplifier le déploiement en production.
\end{enumerate}

\medskip

La démarche adoptée s'organise en phases séquentielles et itératives. Une première
phase d'audit des données sources --- fichiers Excel paroisses, ouvriers et
œuvres, et Shapefile des régions synodales --- permet de valider la structure
réelle des données avant toute modélisation, notamment d'identifier les problèmes
de qualité~: inversion des coordonnées GPS dans les fichiers sources, incohérences
de nommage entre l'Excel et le Shapefile, paroisses sans géolocalisation. Une
deuxième phase de conception définit le schéma de base de données, l'architecture
des services et les modèles UML. L'implémentation procède ensuite couche par
couche~: modèles de données et scripts d'import, API REST avec sécurisation RBAC,
publication GeoServer des couches cartographiques avec styles SLD, développement
du frontend public et de l'interface d'administration. Une validation fonctionnelle
couvre les scénarios d'utilisation des quatre profils d'accès et les tests de
sécurité conformément aux recommandations OWASP~[11].

\medskip

La suite de ce rapport est structurée comme suit. Le \textbf{Chapitre~1} présente
les concepts généraux fondamentaux --- SIG, standards OGC, architecture REST,
contrôle d'accès basé sur les rôles --- et expose l'état de l'art des solutions
de gestion cartographique institutionnelle existantes~; il conclut par un tableau
comparatif et le positionnement de notre solution par rapport à l'existant. Le
\textbf{Chapitre~2} expose l'analyse des besoins à travers un tableau d'exigences
fonctionnelles et non-fonctionnelles ainsi que les diagrammes UML nécessaires à la
compréhension du système, puis présente la conception de l'architecture technique
retenue. Le \textbf{Chapitre~3} détaille l'environnement d'implémentation et les
choix technologiques, présente les résultats à travers les interfaces fonctionnelles
de la plateforme, et propose une discussion critique des forces et des limites de la
solution réalisée. Une \textbf{Conclusion Générale} clôt ce document en synthétisant
les apports du travail, ses limites et les perspectives d'amélioration et de
déploiement.

\clearpage

% ============================================================
%  CHAPITRE 1 — CONCEPTS GÉNÉRAUX ET ÉTAT DE L'ART
% ============================================================
\chapter{Concepts Généraux et État de l'Art}

\begin{chapplan}
  \item \textbf{Section~\ref{sec:concepts} — Concepts Généraux}~: Système
    d'Information Géographique, standards OGC (WMS\,/\,WFS), base de données
    spatiale PostGIS, architecture REST, contrôle d'accès RBAC et conteneurisation Docker.
  \item \textbf{Section~\ref{sec:etat-art} — État de l'Art}~: analyse comparative
    de quatre solutions (GeoEEC~v1.0, ArcGIS~Online, Google~Maps~Platform,
    QGIS\,/\,OpenLayers) selon un format homogène.
  \item \textbf{Section~\ref{sec:bilan} — Bilan et Positionnement}~: tableau
    comparatif multicritère et positionnement de notre solution par rapport à l'existant.
\end{chapplan}

\lettrine[lines=2, loversize=0.05]{C}{e chapitre} pose les fondements théoriques
et contextuels nécessaires à la compréhension du travail réalisé.
La section~\ref{sec:concepts} présente les
concepts généraux mobilisés tout au long de ce rapport~: systèmes d'information
géographique, standards de services cartographiques, base de données spatiale,
architecture REST, contrôle d'accès basé sur les rôles et conteneurisation. La
section~\ref{sec:etat-art} dresse un état de l'art de quatre solutions existantes
répondant partiellement ou totalement aux besoins de l'EEC, analysées selon un
format homogène. La section~\ref{sec:bilan} synthétise cette analyse dans un
tableau comparatif et positionne notre solution par rapport à l'existant.

\section{Concepts Généraux}
\label{sec:concepts}

\subsection{Système d'Information Géographique (SIG)}

Un \kw{Système d'Information Géographique} (\kw{SIG}) est un système informatique
conçu pour la collecte, le stockage, la manipulation, l'analyse et la représentation
de données \kw{géoréférencées}, c'est-à-dire de données auxquelles est associée une
position sur la surface terrestre~[4]. Cette définition, établie
par Longley et al., est aujourd'hui l'une des plus citées dans la littérature
géomatique internationale. Un SIG intègre trois grandes dimensions~: la dimension
\textbf{spatiale} (coordonnées, géométries), la dimension \textbf{attributaire}
(données descriptives associées aux entités géographiques) et la dimension
\textbf{temporelle} (évolution des données dans le temps).

Dans le contexte de ce projet, le SIG constitue la couche centrale permettant de
lier les 858~paroisses importées, 159~districts et 22~régions synodales de l'EEC
à leurs coordonnées géographiques sur le territoire camerounais. Il permet non seulement
la visualisation cartographique de ces entités, mais aussi leur interrogation
spatiale~: identification des paroisses proches d'une localité, calcul de densité
pastorale par région, localisation des zones sous-dotées en ressources. Ces
opérations sont directement utiles au pilotage stratégique de l'institution.

\subsection{Standards OGC et Services Cartographiques Web}

L'\kw{Open Geospatial Consortium} (\kw{OGC}) est un organisme de normalisation
international fondé en 1994, regroupant des industriels, des organismes
gouvernementaux et des institutions académiques du domaine géospatial. Il
publie des \textbf{spécifications ouvertes} pour les services géographiques sur
le web, garantissant l'\textbf{interopérabilité} entre différents serveurs
cartographiques et clients web, indépendamment des fournisseurs et des
plateformes logicielles.

\subsubsection{Web Map Service (WMS)}

Le standard \textit{Web Map Service} (WMS) 1.3.0, publié par l'OGC en 2006,
définit une interface de service permettant à un client web de requêter des
représentations cartographiques sous forme d'images raster depuis un serveur
géographique~[5]. Le client envoie une requête HTTP décrivant
l'étendue géographique souhaitée, les couches à afficher et les dimensions de
l'image~; le serveur retourne une image PNG ou JPEG directement superposable sur
n'importe quelle interface cartographique compatible. Dans ce projet, GeoServer
publie les couches des régions synodales, des paroisses et des œuvres sociales de
l'EEC en WMS~1.3.0, ce qui permet à l'interface Leaflet d'afficher ces couches
institutionnelles par-dessus le fond de carte OpenStreetMap.

\subsubsection{Web Feature Service (WFS)}

Le standard \textit{Web Feature Service} (WFS) 2.0.0, publié par l'OGC en 2010,
permet de requêter et de récupérer des données géographiques vectorielles ---
points, lignes, polygones --- sous forme structurée (GeoJSON ou GML) depuis un
serveur cartographique~[6]. Contrairement au WMS qui transmet des
images prérendues, le WFS retourne des données brutes directement manipulables
côté client. Ce projet exploite le WFS pour permettre à l'interface cartographique
de récupérer les géométries des régions synodales et d'appliquer des filtres
géographiques dynamiques selon le périmètre de l'utilisateur connecté.

\subsection{Base de Données Spatiale PostGIS}

\kw{PostGIS} est une extension spatiale open source pour le système de gestion de
base de données relationnelle \kw{PostgreSQL}~[12]. Elle ajoute la
prise en charge de \textbf{types de données géographiques} conformes aux
spécifications OGC Simple Features --- géométries points (\texttt{PointField}),
lignes (\texttt{LineStringField}) et polygones (\texttt{PolygonField},
\texttt{MultiPolygonField}) --- ainsi que des \textbf{opérateurs spatiaux}
(intersections, calcul de distances, recouvrements) et des fonctions de
transformation de systèmes de coordonnées. PostGIS est le \textbf{standard de
facto} pour les bases de données spatiales open source, utilisé des systèmes
d'information institutionnels aux grandes applications de cartographie
collaborative.

Dans ce projet, PostGIS stocke la géométrie point de chaque paroisse et les
polygones des 22~régions synodales importés depuis le Shapefile source.
GeoDjango, la couche spatiale du framework Django, expose ces données via des
modèles Python et permet l'interrogation spatiale directement depuis l'ORM
(\textit{Object-Relational Mapper}), sans requêtes SQL manuelles.

\subsection{Architecture REST et Interface de Programmation (API)}

Le style architectural \kw{REST} (\textit{Representational State Transfer}) est
défini par Fielding comme un ensemble de contraintes architecturales pour les
systèmes distribués~[13]~: séparation client-serveur, absence
d'état côté serveur (\textit{stateless}), mise en cache des réponses, interface
uniforme et système de couches. Une \kw{API REST} expose des ressources
identifiées par des URI, accessibles via les méthodes HTTP standard
(GET, POST, PUT, PATCH, DELETE), et retourne les données dans un format
standardisé, généralement \textbf{JSON}. Ce style garantit la \textbf{scalabilité},
la \textbf{maintenabilité} et l'\textbf{interopérabilité} des services web.

\kw{Django REST Framework} (\kw{DRF}) est la bibliothèque de référence pour
implémenter des API REST avec le framework Django~[14]. Dans ce projet, la couche
API du backend expose l'ensemble des ressources institutionnelles de l'EEC ---
paroisses, districts, régions, œuvres, ouvriers, statistiques, journaux d'audit
--- consommées par le frontend Next.js~16. Cette séparation stricte entre backend
API et frontend constitue un pilier de l'architecture à services découplés retenue.

\subsection{Contrôle d'Accès Basé sur les Rôles (RBAC)}

Le \kw{contrôle d'accès basé sur les rôles} (\kw{RBAC} --- \textit{Role-Based
Access Control}) est un modèle de sécurité dans lequel les droits d'accès aux
ressources d'un système sont accordés selon le \textbf{rôle} attribué à chaque
utilisateur, et non selon son identité individuelle. Un rôle représente un
ensemble cohérent de \textbf{permissions} correspondant à une fonction ou une
responsabilité dans l'organisation. Ce modèle réduit la complexité
administrative, minimise le risque d'\textbf{escalade de privilèges} et facilite
l'audit de sécurité. L'OWASP le recommande comme mécanisme d'autorisation de
référence pour les applications web~[11].

Dans ce projet, quatre rôles hiérarchiques alignés sur la structure synodale de
l'EEC ont été définis~:
\begin{itemize}[itemsep=2pt, topsep=4pt, leftmargin=1.5cm]
  \item \textbf{SUPER} (Bureau National)~: accès complet à l'ensemble des données
    de toutes les régions synodales~;
  \item \textbf{REGION} (Administrateur régional)~: accès limité aux entités de
    la région synodale dont l'utilisateur est responsable~;
  \item \textbf{DISTRICT} (Administrateur de district)~: accès limité aux entités
    du district ecclésial correspondant~;
  \item \textbf{PAROISSE} (Administrateur paroissial)~: accès en lecture étendue,
    modifications limitées à la paroisse de l'utilisateur.
\end{itemize}

\subsection{Conteneurisation --- Docker et Docker Compose}

La \kw{conteneurisation} est une technique de virtualisation légère qui isole une
application et l'ensemble de ses dépendances dans une unité standardisée appelée
\textbf{conteneur}. Contrairement à la virtualisation classique, un conteneur
partage le noyau du système d'exploitation hôte tout en disposant de son propre
système de fichiers, de ses processus et de sa pile réseau, le rendant bien plus
léger et plus rapide à démarrer qu'une machine virtuelle. \kw{Docker} est la
plateforme de conteneurisation la plus répandue, avec plusieurs millions de
déploiements en production à travers le monde~[15].
\kw{Docker Compose} permet de définir et d'orchestrer des applications
multi-conteneurs dans un fichier de configuration déclaratif
(\texttt{docker-compose.yml}).

Ce projet déploie cinq services orchestrés par Docker Compose~:
\begin{enumerate}[itemsep=2pt, topsep=4pt, leftmargin=1.5cm]
  \item \texttt{backend} --- le serveur Django~5 avec GeoDjango et DRF~;
  \item \texttt{db} --- PostgreSQL~16 avec l'extension PostGIS~3.4~;
  \item \texttt{geoserver} --- GeoServer~2.26 pour la publication WMS/WFS~;
  \item \texttt{redis} --- le serveur de cache Redis~7~;
  \item \texttt{frontend} --- le serveur Next.js~16.
\end{enumerate}

Cette configuration garantit la reproductibilité de l'environnement d'exécution
sur n'importe quelle machine hôte et simplifie considérablement le déploiement
en production, un avantage décisif dans un contexte institutionnel à ressources
informatiques limitées.

% ----
\section{État de l'Art des Solutions Existantes}
\label{sec:etat-art}

Plusieurs catégories de solutions existent pour la gestion cartographique
institutionnelle. Cette section en analyse quatre représentatives~: la plateforme
GeoEEC développée antérieurement, deux plateformes cartographiques commerciales
majeures (ArcGIS Online et Google Maps Platform) et une solution libre de
référence (QGIS associé à OpenLayers). Pour chaque solution, le même format
d'analyse est appliqué~: description générale, illustration, fonctionnalités
principales, forces et faiblesses.

\subsection{Plateforme GeoEEC v1.0 (Version Antérieure)}

La plateforme GeoEEC v1.0 a été développée en 2025 dans le cadre d'un projet
académique antérieur à ce rapport. Elle constitue la référence la plus proche
de notre contexte puisqu'elle traite du même cas d'usage --- la cartographie des
entités de l'EEC --- et a été conçue avec une partie du même stack technologique.
Elle repose sur Django~5 (sans Docker), une authentification JWT
(\textit{JSON Web Token}), un frontend Next.js (Pages Router), une interface
cartographique Leaflet, et un chatbot intégré basé sur le modèle de langage
Ollama phi3:mini associé à LangChain et FAISS. Les données géographiques des
régions synodales sont stockées sous forme de fichiers GeoJSON statiques dans le
répertoire \texttt{public/} du frontend, sans serveur cartographique dédié.

\begin{figure}[H]
  \center
  \includegraphics[width=0.81\textwidth]{ancienne.png}
  \caption{Interface cartographique de la plateforme GeoEEC v1.0}
  \label{fig:diag-contexte}
\end{figure}

\noindent\textbf{Fonctionnalités principales~:}
\begin{itemize}[itemsep=2pt, topsep=4pt, leftmargin=1.5cm]
  \item Carte interactive Leaflet affichant les paroisses, œuvres et ouvriers~;
  \item Import Excel des trois types d'entités avec gestion de la conversion des
    coordonnées sources (Coord\_x/Coord\_y)~;
  \item API REST CRUD pour les paroisses, œuvres et ouvriers~;
  \item Statistiques par région, district et type d'entité~; export PDF~;
  \item Chatbot IA contextuel (Ollama phi3:mini + LangChain + FAISS)~;
  \item Quatre rôles utilisateurs~: admin\_general, admin\_regional,
    utilisateur\_auth, visiteur.
\end{itemize}

\noindent\textbf{Forces~:}
\begin{itemize}[itemsep=2pt, topsep=4pt, leftmargin=1.5cm]
  \item Première implémentation fonctionnelle de la cartographie EEC, validant
    la faisabilité de l'approche Django + Leaflet~;
  \item Module d'import Excel opérationnel avec logique de dépivoter des données
    d'œuvres sociales~;
  \item Chatbot IA innovant permettant des requêtes en langage naturel sur les
    données institutionnelles.
\end{itemize}

\noindent\textbf{Faiblesses et limites~:}
\begin{itemize}[itemsep=2pt, topsep=4pt, leftmargin=1.5cm]
  \item GeoJSON statique~: les couches géographiques des régions synodales ne
    sont pas gérées par un serveur cartographique OGC~; tout changement de
    couche nécessite un redéploiement du frontend~;
  \item Authentification JWT~: les tokens transmis côté client sont plus
    vulnérables aux attaques XSS qu'une authentification par sessions
    \textit{HttpOnly}~;
  \item Absence de journal d'audit~: aucune traçabilité des modifications
    apportées aux données institutionnelles~;
  \item Déploiement non conteneurisé~: dépendances système non isolées,
    portabilité et maintenance complexes~;
  \item RBAC non aligné sur la hiérarchie synodale réelle~: les quatre rôles
    ne correspondent pas aux périmètres district et paroissial de l'EEC~;
  \item Absence de standards OGC~: interopérabilité cartographique limitée, pas
    de WMS ni de WFS exploitables depuis des clients tiers.
\end{itemize}

\subsection{ArcGIS Online (Esri)}

ArcGIS Online est la plateforme SIG en nuage de la société Esri, leader mondial
du logiciel géographique. Elle permet la création, la publication et l'analyse de
cartes interactives depuis un navigateur web, sans installation de logiciel, en
s'appuyant sur une infrastructure cloud hautement disponible. La plateforme intègre
des outils d'analyse spatiale avancée, de gestion de couches de données, de
création de tableaux de bord et de collaboration entre
équipes~[8]. Elle est largement déployée dans des organisations
gouvernementales, des ONG et des institutions académiques.

\begin{figure}[H]
  \center
  \includegraphics[width=0.81\textwidth]{aris.png}
  \caption{lateforme cartographique ArcGIS Online (Esri, 2024}
  \label{fig:diag-contexte}
\end{figure}

\noindent\textbf{Fonctionnalités principales~:}
\begin{itemize}[itemsep=2pt, topsep=4pt, leftmargin=1.5cm]
  \item Carte web interactive avec accès à de nombreux fonds de carte et couches
    de données mondiales~;
  \item Outils d'analyse spatiale avancés~: zones tampons, interpolations, analyses
    de réseaux, clustering~;
  \item Tableau de bord configurable avec indicateurs actualisables en temps réel~;
  \item Contrôle des permissions sur les couches (lecture, édition, partage)~;
  \item Applications mobiles intégrées pour la collecte de données terrain~;
  \item APIs REST compatibles avec les standards OGC WMS et WFS.
\end{itemize}

\noindent\textbf{Forces~:}
\begin{itemize}[itemsep=2pt, topsep=4pt, leftmargin=1.5cm]
  \item Plateforme mature avec un très large éventail de fonctionnalités SIG~;
  \item Interface utilisateur intuitive, accessible aux utilisateurs non
    techniciens~;
  \item Infrastructure cloud scalable avec haute disponibilité et support
    professionnel.
\end{itemize}

\noindent\textbf{Faiblesses et limites~:}
\begin{itemize}[itemsep=2pt, topsep=4pt, leftmargin=1.5cm]
  \item Licence commerciale onéreuse~: les abonnements annuels ArcGIS Online sont
    prohibitifs pour une institution confessionnelle à ressources limitées~;
  \item Aucune logique hiérarchique institutionnelle configurable~: pas de notion
    de périmètre synodal, de district ou de paroisse dans le modèle d'accès~;
  \item Données hébergées sur les serveurs d'Esri (États-Unis)~: dépendance à une
    infrastructure cloud étrangère, souveraineté des données compromise~;
  \item Absence de gestion intégrée des données métier EEC~: ouvriers
    ecclésiastiques, œuvres sociales et statistiques annuelles ne peuvent être
    gérés nativement.
\end{itemize}

\subsection{Google Maps Platform}

Google Maps Platform est l'ensemble des APIs cartographiques de Google, permettant
d'intégrer des cartes interactives, des fonctionnalités de géocodage et des
itinéraires dans des applications web et mobiles~[9]. C'est la
solution cartographique web la plus utilisée mondialement. Elle s'appuie sur le
fond de carte Google, qui bénéficie d'une couverture très détaillée, y compris
pour les villes et localités camerounaises.

\begin{figure}[H]
  \center
  \includegraphics[width=0.81\textwidth]{google.png}
  \caption{Google Maps Platform }
  \label{fig:diag-contexte}
\end{figure}
\noindent\textbf{Fonctionnalités principales~:}
\begin{itemize}[itemsep=2pt, topsep=4pt, leftmargin=1.5cm]
  \item Cartographie interactive avec fond de carte Google haute résolution~;
  \item Géocodage (conversion adresse vers coordonnées GPS et inversement)~;
  \item Calcul d'itinéraires multimodaux~;
  \item Marqueurs personnalisés et couches de données~;
  \item APIs REST bien documentées, intégration simple pour les développeurs~;
  \item SDK mobiles Android et iOS.
\end{itemize}

\noindent\textbf{Forces~:}
\begin{itemize}[itemsep=2pt, topsep=4pt, leftmargin=1.5cm]
  \item Couverture cartographique mondiale très détaillée, performante pour le
    Cameroun~;
  \item API simple à intégrer, documentation exhaustive, large communauté de
    développeurs~;
  \item Géocodage très performant pour localiser des adresses camerounaises.
\end{itemize}

\noindent\textbf{Faiblesses et limites~:}
\begin{itemize}[itemsep=2pt, topsep=4pt, leftmargin=1.5cm]
  \item Modèle de facturation à l'usage~: au-delà du quota gratuit mensuel, les
    coûts peuvent devenir significatifs pour un usage institutionnel~;
  \item Pas de serveur cartographique OGC~: impossible de publier des couches
    propres en WMS ou WFS~;
  \item Aucune gestion des entités institutionnelles ni de la hiérarchie
    synodalo-ecclésiale~;
  \item Données cartographiques propriétaires hébergées sur les serveurs de
    Google~: dépendance totale et absence de souveraineté~;
  \item Pas de gestion des rôles, de journal d'audit ni d'import de données
    institutionnelles~;
  \item Conformité aux obligations légales locales et souveraineté des données
    problématiques pour une institution camerounaise.
\end{itemize}

\subsection{QGIS et OpenLayers}

QGIS est un Système d'Information Géographique libre et open source parmi les
plus utilisés au monde dans les secteurs académique, institutionnel et de la
coopération internationale~[10]. Il offre des fonctionnalités
complètes d'édition, d'analyse et de visualisation cartographique depuis une
application bureau. OpenLayers est une bibliothèque JavaScript open source
permettant d'afficher des cartes interactives dans un navigateur web, avec support
natif des standards OGC WMS et WFS.

\begin{figure}[H]
  \center
  \includegraphics[width=0.81\textwidth]{qgis.png}
  
  \caption{Interface du logiciel SIG bureau QGIS }
  \label{fig:diag-contexte}
\end{figure}

\noindent\textbf{Fonctionnalités principales~:}
\begin{itemize}[itemsep=2pt, topsep=4pt, leftmargin=1.5cm]
  \item Édition et analyse cartographique complète (QGIS)~;
  \item Symbologie avancée et mise en page de cartes imprimées~;
  \item Traitement géospatial avancé~: zones tampons, interpolations, analyses de
    réseau, calculs rasters~;
  \item Affichage de couches WMS et WFS dans le navigateur (OpenLayers)~;
  \item Support natif des formats SHP, GeoJSON, PostGIS et GeoTIFF.
\end{itemize}

\noindent\textbf{Forces~:}
\begin{itemize}[itemsep=2pt, topsep=4pt, leftmargin=1.5cm]
  \item Open source, sans coût de licence, avec une large communauté active et
    une documentation exhaustive~;
  \item Fonctionnalités SIG professionnelles complètes~;
  \item Haute interopérabilité grâce au support natif des standards OGC.
\end{itemize}

\noindent\textbf{Faiblesses et limites~:}
\begin{itemize}[itemsep=2pt, topsep=4pt, leftmargin=1.5cm]
  \item QGIS est une application bureau~: son déploiement en usage institutionnel
    web partagé est impraticable~;
  \item Aucun contrôle d'accès multi-utilisateurs basé sur des périmètres
    institutionnels hiérarchisés~;
  \item OpenLayers seul ne fournit ni backend, ni gestion de rôles, ni import de
    données institutionnelles, ni statistiques~;
  \item Nécessite des compétences SIG spécialisées~: inadapté aux administrateurs
    régionaux et paroissiaux non techniciens de l'EEC~;
  \item Déploiement complexe en environnement institutionnel à infrastructure
    informatique limitée.
\end{itemize}

% ----
\section{Bilan et Positionnement}
\label{sec:bilan}

\subsection{Positionnement de Notre Solution}

Cette sous-section positionne qualitativement notre solution par rapport à
l'existant~; une comparaison chiffrée détaillée (tableau~\ref{tab:comparatif})
est présentée en section~\ref{ssec:disc-comparatif} du Chapitre~3, une fois les
résultats concrets de l'implémentation exposés.

L'analyse comparative met en évidence que notre solution est la seule à réunir
simultanément l'ensemble des caractéristiques requises pour une plateforme
institutionnelle cartographique au service de l'EEC.

Les plateformes commerciales --- ArcGIS Online et Google Maps Platform --- offrent
des capacités cartographiques puissantes mais ne répondent ni aux contraintes
économiques de l'institution ni à ses exigences de hiérarchie institutionnelle.
Leurs modèles de licences commerciales et leur dépendance à des infrastructures
cloud étrangères les rendent inadaptées à une institution confessionnelle opérant
au Cameroun. Elles ne proposent aucun mécanisme de gestion des données métier
spécifiques à l'EEC --- ouvriers ecclésiastiques, œuvres sociales, statistiques
annuelles des paroisses --- ni de contrôle d'accès structuré selon les périmètres
synodaux.

QGIS et OpenLayers, bien que techniquement solides et libres de droits, ne
constituent pas une plateforme web institutionnelle multi-utilisateurs~: QGIS est
une application bureau inadaptée aux administrateurs régionaux non techniciens, et
OpenLayers est une bibliothèque de rendu client sans backend, sans gestion de
données et sans contrôle d'accès.

La plateforme GeoEEC v1.0 constitue le précédent le plus proche de notre contexte,
mais présente des lacunes critiques identifiées lors de l'audit technique~:
vulnérabilité de l'authentification JWT par rapport aux sessions HttpOnly, absence
de journal d'audit, couches géographiques statiques non publiées via un serveur
OGC, RBAC non aligné sur la hiérarchie synodale réelle à quatre niveaux, et absence
de conteneurisation.

Notre solution apporte une valeur ajoutée précise et documentée sur chacun de ces
axes. Elle est la première à implémenter simultanément~: une architecture de services
découplés et entièrement conteneurisée, un serveur cartographique OGC natif
(GeoServer~2.26) publiant des couches WMS et WFS interopérables, une base de données
spatiale PostGIS avec gestion des géométries institutionnelles, une authentification
sécurisée par sessions Django (\textit{HttpOnly}, \textit{Secure},
\textit{SameSite=Strict}) et un RBAC exactement aligné sur les quatre niveaux
hiérarchiques de l'EEC, avec un journal d'audit complet et automatique de toutes
les opérations d'écriture sur les données.

\medskip
\noindent\textbf{Bilan du chapitre.}
Ce chapitre a établi les fondements théoriques du travail réalisé et mis en
évidence, par une analyse comparative rigoureuse, la pertinence et l'originalité
de l'approche proposée. Les six concepts fondamentaux --- SIG, standards OGC
(WMS/WFS), PostGIS, architecture REST, RBAC et Docker --- constituent le socle sur
lequel repose l'architecture décrite dans les chapitres suivants. L'état de l'art
de quatre solutions existantes a démontré qu'aucune d'elles ne couvre simultanément
l'ensemble des exigences institutionnelles de l'EEC. Le \textbf{Chapitre~2} procédera
à l'analyse détaillée des besoins du système à travers les exigences fonctionnelles
et non-fonctionnelles et les modélisations UML, puis présentera les décisions de
conception de l'architecture technique retenue.

\clearpage

% ============================================================
%  CHAPITRE 2 — ANALYSE ET CONCEPTION
% ============================================================
\chapter{Analyse et Conception du Système de Géolocalisation EEC}
\label{chap:analyse-conception}

\begin{chapplan}
  \item \textbf{Section~\ref{sec:audit-donnees} --- Audit et Préparation des Données}~:
    \begin{itemize}[noitemsep, topsep=3pt, leftmargin=1.2cm]
      \item fichiers sources analysés~: paroisses, catégorisation officielle, ouvriers, œuvres~;
      \item anomalies détectées~: doublons, incohérences, coordonnées inversées, structures hétérogènes~;
      \item méthodologie de nettoyage et de rapprochement flou outillée par scripts Python~;
      \item résultats chiffrés~: 858~paroisses intégrées, 333~correspondances établies.
    \end{itemize}
  \item \textbf{Section~\ref{sec:analyse} --- Analyse des Besoins}~:
    \begin{itemize}[noitemsep, topsep=3pt, leftmargin=1.2cm]
      \item ingénierie des exigences~: 24~exigences fonctionnelles et 12~non-fonctionnelles~;
      \item diagramme de contexte~: six acteurs externes identifiés~;
      \item diagramme de packages~: trois couches (frontend, backend, infrastructure)~;
      \item cas d'utilisation UC01--UC29 avec tableau de description structuré~;
      \item diagramme de classes métier~: huit entités du domaine institutionnel~;
      \item diagramme d'activité~: flux d'import de données depuis Excel.
    \end{itemize}
  \item \textbf{Section~\ref{sec:conception} --- Conception du Système}~:
    \begin{itemize}[noitemsep, topsep=3pt, leftmargin=1.2cm]
      \item architecture technique~: cinq services Docker découplés et leurs communications~;
      \item modèle de données~: quinze tables PostGIS avec champs géospatiaux~;
      \item diagramme de séquence~: flux d'authentification sécurisée par sessions Django~;
      \item diagramme de classes technique~: modèles Django et relations inter-classes.
    \end{itemize}
  \item \textbf{Section~\ref{sec:modelisation-math} --- Modélisation Mathématique}~:
    \begin{itemize}[noitemsep, topsep=3pt, leftmargin=1.2cm]
      \item modélisation du réseau routier comme graphe pondéré $G=(V,E,w)$~;
      \item formule de Haversine pour la distance à vol d'oiseau~;
      \item algorithmes Dijkstra et A\* pour le calcul de plus court chemin.
    \end{itemize}
\end{chapplan}

\lettrine[lines=2, loversize=0.05]{L}{a} compréhension précise des besoins et
la définition rigoureuse d'une architecture sont les deux prérequis fondamentaux
d'une implémentation réussie. Ce chapitre s'ouvre par la
section~\ref{sec:audit-donnees}, consacrée à l'audit et à la préparation des
données sources --- étape préalable à toute modélisation, tant les
incohérences relevées dans les fichiers fournis par l'EEC ont conditionné
les choix de conception qui suivent. Ce chapitre traduit ensuite les
contraintes institutionnelles et techniques identifiées au Chapitre~1 en
exigences mesurables et en décisions de conception concrètes. La
section~\ref{sec:analyse} formalise les besoins du
système à travers un tableau d'exigences fonctionnelles et non-fonctionnelles,
puis une série de modèles UML couvrant les acteurs, les modules, les cas
d'utilisation, les entités du domaine et les flux de traitement.
La section~\ref{sec:conception} s'appuie sur cette analyse pour définir
l'architecture technique retenue, le schéma de la base de données spatiale,
les diagrammes de séquence et de classes de l'implémentation. Enfin, la
section~\ref{sec:modelisation-math} formalise mathématiquement le service de
calcul d'itinéraire~: modélisation du réseau routier en graphe pondéré,
formule de Haversine et algorithmes de recherche de plus court chemin
(Dijkstra, A\*).

%============================================================
\section{Audit et Préparation des Données}
\label{sec:audit-donnees}
%============================================================

\subsection{Contexte et Objectifs de l'Audit}

Avant toute modélisation, une phase d'audit approfondi des données sources
fournies par l'EEC a été menée. Cette étape s'imposait pour une raison
simple~: la qualité d'une plateforme de géolocalisation dépend directement
de la qualité des données qu'elle expose~; modéliser un schéma de base de
données sans connaître précisément la structure réelle, les incohérences et
les lacunes des fichiers sources aurait exposé le projet à des
retraitements coûteux en cours d'implémentation. L'audit visait trois
objectifs~: (1)~inventorier exhaustivement les fichiers disponibles et leur
structure réelle~; (2)~détecter et catégoriser les anomalies (doublons,
incohérences, données manquantes)~; (3)~définir une méthodologie de
nettoyage et de rapprochement reproductible, outillée par des scripts
Python dédiés plutôt que par des corrections manuelles ponctuelles.

\subsection{Fichiers de Données Analysés}

Quatre familles de fichiers hétérogènes ont été soumises à l'audit~:
\begin{itemize}[itemsep=2pt, topsep=4pt, leftmargin=1.5cm]
  \item un tableur \textbf{paroisses} issu d'une enquête de géolocalisation
    de terrain, comportant les coordonnées GPS relevées et une structure en
    deux feuilles (données brutes et données consolidées)~;
  \item un document \textbf{officiel de catégorisation} des paroisses
    (format traitement de texte), listant le nom, la région synodale et la
    classe (catégorie C1--C4, B1--B2, A1--A2, A++) de chaque paroisse
    reconnue par le Bureau National~;
  \item un tableur \textbf{ouvriers ecclésiastiques}, où chaque ligne
    combine une paroisse et la liste de ses ouvriers dans une colonne texte
    libre (\og Grade \fg{}), plutôt que des enregistrements individuels~;
  \item un tableur \textbf{œuvres sociales et éducatives}, structuré en
    tableau croisé (\textit{pivot})~: une ligne par région synodale, les
    types d'œuvres en colonnes --- format inexploitable directement par un
    modèle relationnel, qui suppose une ligne par entité.
\end{itemize}

\subsection{Anomalies et Difficultés Rencontrées}

L'audit a mis en évidence des anomalies de nature et de gravité variables,
résumées au tableau~\ref{tab:audit-anomalies}.

\begin{table}[H]
\centering
\caption{Principales anomalies détectées lors de l'audit des données}
\label{tab:audit-anomalies}
\setlength{\tabcolsep}{5pt}
\renewcommand{\arraystretch}{1.22}
\resizebox{\linewidth}{!}{%
\begin{tabular}{>{\bfseries}m{3.6cm} m{8.2cm} m{2.4cm}}
\toprule
\rowcolor{eecgreen}
\multicolumn{1}{>{\bfseries\color{white}}m{3.6cm}}{Anomalie} &
\multicolumn{1}{>{\bfseries\color{white}}m{8.2cm}}{Description} &
\multicolumn{1}{>{\bfseries\color{white}}m{2.4cm}}{Sévérité} \\
\midrule
\rowcolor{eeclightgreen!60}
Doublons de paroisses & Une même paroisse enregistrée plusieurs fois sous des
  intitulés légèrement différents (ex.~\og Jourdain de Beka Hossere \fg{} et
  \og Le Jourdain de Beka Hossere \fg{}) & Haute \\
Incohérences de nommage & Écarts orthographiques entre le fichier d'enquête
  et le fichier officiel de catégorisation (ex.~\og Bamemba \fg{} vs
  \og Bameba \fg{}), ainsi qu'entre les noms de régions synodales du tableur
  et du Shapefile & Haute \\
\rowcolor{eeclightgreen!60}
Coordonnées GPS inversées & Dans le fichier source, la colonne
  \texttt{Coord\_x} contient en réalité la latitude et \texttt{Coord\_y} la
  longitude --- l'inverse de la convention géographique standard & Haute \\
Données manquantes & Absence de coordonnées GPS pour une majorité de
  paroisses~; absence de classe (catégorie) pour les paroisses non recensées
  dans le fichier officiel & Haute \\
\rowcolor{eeclightgreen!60}
Structure hétérogène & Le fichier des œuvres sociales est organisé en
  tableau croisé (région en ligne, type d'œuvre en colonne) et nécessite une
  opération de \og dépivotage \fg{} avant tout import relationnel & Moyenne \\
Enregistrements composites & Le fichier des ouvriers combine plusieurs
  ouvriers par paroisse dans une seule cellule texte, empêchant un import
  direct en enregistrements individuels & Moyenne \\
\bottomrule
\end{tabular}}
\end{table}

Le rapprochement entre les deux sources de vérité relatives aux
paroisses --- l'enquête de géolocalisation et le fichier officiel de
catégorisation --- a constitué la difficulté la plus significative de
l'audit~: les deux fichiers, produits indépendamment à des dates et par des
services différents, ne partageaient ni identifiant commun ni orthographe
strictement normalisée pour les noms de paroisses.

\subsection{Méthodologie de Nettoyage et de Rapprochement}

Une méthodologie outillée, plutôt qu'une correction manuelle fichier par
fichier, a été retenue afin de garantir la reproductibilité du traitement à
chaque nouvelle livraison de données par l'EEC. Elle s'appuie sur une suite
de scripts Python dédiés (répertoire \texttt{backend/data\_audit/})~:
\begin{enumerate}[itemsep=2pt, topsep=4pt, leftmargin=1.5cm]
  \item \textbf{Inventaire structurel} --- inspection automatique de chaque
    fichier source (types de colonnes, taux de remplissage, valeurs
    distinctes) afin d'établir un état des lieux exhaustif avant tout
    traitement~;
  \item \textbf{Extraction du référentiel officiel} --- conversion du
    document de catégorisation (format traitement de texte) en table
    structurée exploitable par les scripts suivants~;
  \item \textbf{Nettoyage et normalisation} --- correction de la casse, des
    espaces superflus, harmonisation des intitulés de région synodale entre
    le tableur d'enquête et le Shapefile, inversion corrective des colonnes
    de coordonnées GPS~;
  \item \textbf{Rapprochement flou} (\textit{fuzzy matching}) --- comparaison
    algorithmique de chaque nom de paroisse de l'enquête avec le référentiel
    officiel par un score de similarité textuelle, classant chaque
    correspondance en \textit{exacte}, \textit{approchée} (au-delà d'un
    seuil de similarité) ou \textit{absente}~;
  \item \textbf{Intégration des classes} --- attribution de la catégorie
    officielle (C1--C4, B1--B2, A1--A2, A++) à chaque paroisse rapprochée
    avec succès, et création des paroisses présentes uniquement dans le
    référentiel officiel (sans coordonnées GPS, district à préciser)~;
  \item \textbf{Génération d'un rapport de comparaison} --- production d'un
    document de traçabilité exhaustif, ligne par ligne, de chaque décision
    de rapprochement, soumis à validation avant intégration définitive en
    base.
\end{enumerate}

\subsection{Résultats de l'Audit}

Le tableau~\ref{tab:audit-resultats} synthétise les résultats obtenus après
application de cette méthodologie sur le jeu de données des paroisses.

\begin{table}[H]
\centering
\caption{Résultats du rapprochement des données de paroisses}
\label{tab:audit-resultats}
\setlength{\tabcolsep}{5pt}
\renewcommand{\arraystretch}{1.22}
\begin{tabular}{>{\bfseries}m{7.5cm} >{\centering\arraybackslash}m{3.5cm}}
\toprule
\rowcolor{eecgreen}
\multicolumn{1}{>{\bfseries\color{white}}m{7.5cm}}{Indicateur} &
\multicolumn{1}{>{\bfseries\color{white}\centering\arraybackslash}m{3.5cm}}{Valeur} \\
\midrule
\rowcolor{eeclightgreen!60}
Paroisses issues de l'enquête de géolocalisation & 516 \\
Paroisses ajoutées depuis le référentiel officiel (sans GPS) & 354 \\
\rowcolor{eeclightgreen!60}
Total de paroisses intégrées en base & 858 \\
Correspondances exactes (nom identique) & 267 \\
\rowcolor{eeclightgreen!60}
Correspondances approchées (similarité orthographique) & 66 \\
Paroisses sans correspondance officielle (classe non définie) & 183 \\
\rowcolor{eeclightgreen!60}
Paroisses géolocalisées (coordonnées GPS valides) & 378 \\
Paroisses sans coordonnées GPS valides & 480 \\
\bottomrule
\end{tabular}
\end{table}

Ces résultats appellent deux commentaires. Premièrement, le taux de
correspondance automatique (267~exactes~+~66~approchées, soit
333~paroisses sur 516 issues de l'enquête, environ 65~\%) confirme que le
rapprochement flou était nécessaire~: une comparaison stricte des chaînes de
caractères aurait manqué les 66~correspondances approchées, soit près de
13~\% de rapprochements supplémentaires. Deuxièmement, les
183~paroisses sans correspondance officielle et les 480~paroisses sans
coordonnées GPS ne sont pas silencieusement écartées~: elles sont intégrées
en base avec un statut explicite (\texttt{categorie} vide, \texttt{position}
nulle), consultables et corrigibles depuis l'interface d'administration,
plutôt que d'être perdues ou approximées automatiquement. C'est directement
ce constat --- la nécessité d'un contrôle humain sur les cas ambigus --- qui a
motivé le retrait du module d'import automatisé de masse
(section~\ref{ssec:activite}) au profit d'une saisie et d'une correction
supervisées dans l'application elle-même.

%------------------------------------------------------------
\section{Analyse des Besoins}
\label{sec:analyse}
%------------------------------------------------------------

\subsection{Ingénierie des Exigences}
\label{ssec:exigences}

L'ingénierie des exigences consiste à identifier, formaliser et prioriser les
propriétés que le système doit satisfaire avant toute activité de conception.
Deux catégories sont distinguées~: les \textbf{exigences fonctionnelles}
(\textbf{EF}), qui décrivent ce que le système doit réaliser, et les
\textbf{exigences non-fonctionnelles} (\textbf{ENF}), qui portent sur les
contraintes qualitatives de fonctionnement (sécurité, performance, portabilité,
interopérabilité). Ces exigences sont directement dérivées des six objectifs
formulés en Introduction Générale et des lacunes identifiées dans l'état de
l'art du Chapitre~1.

\subsubsection{Exigences Fonctionnelles}

Le Tableau~\ref{tab:ef} recense les vingt-quatre exigences fonctionnelles du
système dans sa version actuelle, organisées par domaine~: authentification
et contrôle d'accès, cartographie, administration des entités, export,
statistiques, traçabilité, gestion hiérarchique des comptes et services du
visiteur authentifié (EF18--EF24), dont le calcul d'itinéraire réel et la
navigation guidée en trois dimensions. Conformément à l'évolution du projet
documentée dans l'audit technique (section~\ref{sec:audit-donnees}), le
module d'import de fichiers Excel, présent dans une version antérieure de la
plateforme, a été retiré du périmètre fonctionnel~; seul l'export subsiste.

\setlength{\tabcolsep}{5pt}
\renewcommand{\arraystretch}{1.22}
\begin{longtable}{%
  >{\bfseries\color{eccdarkgreen}}m{1.2cm}
  m{7.8cm}
  m{3.5cm}
  >{\centering\arraybackslash\bfseries}m{1.7cm}}
\caption{Exigences fonctionnelles du système}\label{tab:ef} \\
\toprule
\rowcolor{eecgreen}
\multicolumn{1}{>{\bfseries\color{white}}m{1.2cm}}{ID} &
\multicolumn{1}{>{\bfseries\color{white}}m{7.8cm}}{Description} &
\multicolumn{1}{>{\bfseries\color{white}}m{3.5cm}}{Objectif} &
\multicolumn{1}{>{\bfseries\color{white}\centering\arraybackslash}m{1.7cm}}{Priorité} \\
\midrule
\endfirsthead
\caption[]{\textit{(suite)~Exigences fonctionnelles du système}} \\
\toprule
\rowcolor{eecgreen}
\multicolumn{1}{>{\bfseries\color{white}}m{1.2cm}}{ID} &
\multicolumn{1}{>{\bfseries\color{white}}m{7.8cm}}{Description} &
\multicolumn{1}{>{\bfseries\color{white}}m{3.5cm}}{Objectif} &
\multicolumn{1}{>{\bfseries\color{white}\centering\arraybackslash}m{1.7cm}}{Priorité} \\
\midrule
\endhead
\midrule
\multicolumn{4}{r}{\small\textit{ }} \\
\endfoot
\bottomrule
\endlastfoot
\rowcolor{eeclightgreen!60}
EF01 & Le système doit permettre à tout utilisateur d'accéder à un formulaire
  de connexion sécurisé et de se déconnecter en invalidant la session active. &
  Sécurité & Haute \\
EF02 & Le système doit gérer quatre rôles hiérarchiques administratifs
  --- SUPER, REGION, DISTRICT, PAROISSE --- chacun disposant d'un périmètre
  d'accès aux données strictement défini, ainsi qu'un cinquième profil
  VISITEUR non hiérarchique réservé à l'accès public et personnel. &
  RBAC synodal & Haute \\
\rowcolor{eeclightgreen!60}
EF03 & Le système doit afficher une carte interactive Leaflet.js superposant
  les couches institutionnelles publiées par GeoServer~: paroisses (points),
  régions synodales (polygones) et œuvres (points). & Cartographie & Haute \\
EF04 & La carte doit regrouper automatiquement les marqueurs proches en clusters
  numériques, révélant les entités individuelles à mesure que l'utilisateur
  zoome. & Cartographie & Haute \\
\rowcolor{eeclightgreen!60}
EF05 & La carte doit proposer des filtres hiérarchiques multicritères~:
  région synodale, district, type d'entité (paroisse, œuvre, ouvrier) et
  classe de paroisse (catégorie C1--C4, B1--B2, A1--A2, A++). &
  Cartographie & Haute \\
EF06 & Un clic sur un marqueur de la carte doit ouvrir une popup affichant
  les informations détaillées de l'entité (nom, type, contacts, statut GPS). &
  Cartographie & Haute \\
\rowcolor{eeclightgreen!60}
EF07 & Les administrateurs autorisés doivent pouvoir créer, lire, modifier et
  supprimer des paroisses dans les limites de leur périmètre synodal, y
  compris l'attribution de leur classe (catégorie). &
  Administration & Haute \\
EF08 & Les administrateurs autorisés doivent pouvoir créer, lire, modifier et
  supprimer des œuvres sociales et éducatives dans leur périmètre. &
  Administration & Haute \\
\rowcolor{eeclightgreen!60}
EF09 & Les administrateurs autorisés doivent pouvoir créer, lire, modifier et
  supprimer des ouvriers ecclésiastiques et gérer leurs affectations à une
  paroisse. & Administration & Haute \\
EF10 & Le système doit permettre l'export des listes d'entités au format
  Excel (.xlsx) et PDF depuis l'interface d'administration. & Export & Moyenne \\
\rowcolor{eeclightgreen!60}
EF11 & Le tableau de bord doit afficher des indicateurs synthétiques~:
  nombre de paroisses, districts, ouvriers, œuvres et taux de couverture
  GPS par région. & Statistiques & Haute \\
EF12 & Le système doit permettre la saisie, la modification et la
  consultation des statistiques annuelles de chaque paroisse (membres,
  baptisés, enfants, finances). & Statistiques & Haute \\
\rowcolor{eeclightgreen!60}
EF13 & Toute opération d'écriture (création, modification, suppression) sur
  les données institutionnelles doit être enregistrée automatiquement dans
  un journal d'audit horodaté. & Traçabilité & Haute \\
EF14 & Le journal d'audit doit être consultable par le rôle SUPER, avec des
  filtres par type d'action, entité concernée et période temporelle. &
  Traçabilité & Moyenne \\
\rowcolor{eeclightgreen!60}
EF15 & Le système doit proposer une recherche d'entités (paroisses, œuvres,
  ouvriers) par nom, combinable avec les filtres de région, district et
  classe de paroisse, depuis la carte et depuis l'administration. &
  Ergonomie & Moyenne \\
EF16 & Les visiteurs non authentifiés doivent pouvoir consulter la carte
  publique et les informations générales des paroisses sans accéder au module
  d'administration. & Accès public & Moyenne \\
\rowcolor{eeclightgreen!60}
EF17 & Le rôle SUPER doit pouvoir créer, désactiver et paramétrer tout
  compte administrateur~; les rôles REGION et DISTRICT doivent pouvoir
  créer des comptes administrateurs limités à leur propre périmètre
  synodal (délégation hiérarchique de la gestion des comptes). &
  Gestion comptes & Haute \\
EF18 & Le système doit permettre à un visiteur de s'inscrire avec un email
  et un mot de passe~; un profil visiteur est créé automatiquement à
  l'inscription et persiste entre les sessions. &
  Visiteur auth. & Haute \\
\rowcolor{eeclightgreen!60}
EF19 & Le visiteur authentifié doit pouvoir enregistrer des paroisses en
  favoris depuis la carte et les retrouver dans son espace personnel. &
  Visiteur auth. & Haute \\
EF20 & Le système doit proposer deux modes d'affichage cartographique~: la
  carte standard en deux dimensions (Leaflet) pour la consultation courante,
  et une vue de navigation guidée en trois dimensions (MapLibre~GL) pendant
  un déplacement actif. & Cartographie & Haute \\
\rowcolor{eeclightgreen!60}
EF21 & Le système doit calculer un itinéraire réel sur le réseau routier
  entre un point de départ et une entité de destination, pour trois modes
  de déplacement (voiture, vélo, à pied), avec distance, durée et tracé
  détaillé~; en cas d'indisponibilité du service de routing, un repli sur
  la distance à vol d'oiseau (Haversine) est appliqué automatiquement. &
  Visiteur auth. & Haute \\
EF22 & Une fois l'itinéraire calculé, le visiteur doit pouvoir démarrer une
  navigation guidée en~3D avec suivi de sa position GPS en temps réel,
  recalage sur le tracé et détection automatique de l'arrivée. &
  Visiteur auth. & Haute \\
\rowcolor{eeclightgreen!60}
EF23 & Le système doit conserver l'historique des recherches géographiques
  du visiteur authentifié et lui permettre de le consulter et de le vider
  depuis son profil personnel. &
  Visiteur auth. & Moyenne \\
EF24 & Chaque profil doit être restreint aux seules actions et données de
  son périmètre~: un visiteur non authentifié n'a accès ni aux favoris, ni
  à l'itinéraire, ni à l'historique~; un administrateur REGION, DISTRICT ou
  PAROISSE n'agit que sur les entités de son périmètre synodal, quelle que
  soit l'opération demandée. &
  Sécurité & Haute \\
\end{longtable}

\subsubsection{Exigences Non-Fonctionnelles}

Le Tableau~\ref{tab:enf} liste les douze exigences non-fonctionnelles
définissant les contraintes qualitatives de la plateforme sur les axes
sécurité, performance, interopérabilité, portabilité, maintenabilité,
fiabilité et performance temps réel (ces deux dernières, ENF11 et ENF12,
ajoutées pour couvrir la dépendance aux services de routing externes et les
exigences de fluidité de la navigation~3D).

\setlength{\tabcolsep}{5pt}
\renewcommand{\arraystretch}{1.22}
\begin{longtable}{%
  >{\bfseries\color{eccdarkgreen}}m{1.2cm}
  >{\bfseries}m{2.8cm}
  m{8.2cm}
  >{\centering\arraybackslash\bfseries}m{1.7cm}}
\caption{Exigences non-fonctionnelles du système}\label{tab:enf} \\
\toprule
\rowcolor{eecgreen}
\multicolumn{1}{>{\bfseries\color{white}}m{1.2cm}}{ID} &
\multicolumn{1}{>{\bfseries\color{white}}m{2.8cm}}{Domaine} &
\multicolumn{1}{>{\bfseries\color{white}}m{8.2cm}}{Description} &
\multicolumn{1}{>{\bfseries\color{white}\centering\arraybackslash}m{1.7cm}}{Priorité} \\
\midrule
\endfirsthead
\caption[]{\textit{(suite)~Exigences non-fonctionnelles du système}} \\
\toprule
\rowcolor{eecgreen}
\multicolumn{1}{>{\bfseries\color{white}}m{1.2cm}}{ID} &
\multicolumn{1}{>{\bfseries\color{white}}m{2.8cm}}{Domaine} &
\multicolumn{1}{>{\bfseries\color{white}}m{8.2cm}}{Description} &
\multicolumn{1}{>{\bfseries\color{white}\centering\arraybackslash}m{1.7cm}}{Priorité} \\
\midrule
\endhead
\midrule
\multicolumn{4}{r}{\small\textit{(suite page suivante)}} \\
\endfoot
\bottomrule
\endlastfoot
\rowcolor{eeclightgreen!60}
ENF01 & Sécurité & Les sessions sont gérées par des cookies Django
  \textit{HttpOnly}, \textit{Secure} et \textit{SameSite=Strict}. Chaque
  requête modifiante est protégée par un token CSRF. Les mots de passe sont
  hachés avec PBKDF2-SHA256. & Haute \\
ENF02 & Performance & Le temps de réponse des requêtes API standard (liste,
  détail) est inférieur à 2~secondes pour les jeux de données courants du
  système. & Haute \\
\rowcolor{eeclightgreen!60}
ENF03 & Disponibilité & Le déploiement Docker Compose sur un serveur
  institutionnel assure la continuité des services sans dépendance à une
  infrastructure cloud tierce. & Haute \\
ENF04 & Maintenabilité & L'architecture découplée (backend Django/DRF séparé
  du frontend Next.js) permet la mise à jour indépendante de chaque couche
  sans interruption des autres services. & Haute \\
\rowcolor{eeclightgreen!60}
ENF05 & Interopérabilité & Les couches géographiques sont publiées selon les
  standards OGC WMS~1.3.0 et WFS~2.0.0, garantissant leur consommation par
  tout client cartographique compatible. & Haute \\
ENF06 & Portabilité & La plateforme complète (cinq services) est conteneurisée
  avec Docker et Docker Compose, permettant le déploiement sur tout serveur
  Linux disposant de Docker. & Haute \\
\rowcolor{eeclightgreen!60}
ENF07 & Intégrité & Les formulaires de saisie déclenchent une validation
  systématique des coordonnées GPS (plage du territoire camerounais), du
  nommage des entités, de la classe de paroisse et des relations
  hiérarchiques avant tout enregistrement. & Haute \\
ENF08 & Traçabilité & Chaque entrée d'audit contient~: l'auteur, le type
  d'action (CREATE / UPDATE / DELETE), l'entité cible, l'identifiant de
  l'objet et l'horodatage UTC. & Haute \\
\rowcolor{eeclightgreen!60}
ENF09 & Ergonomie & L'interface d'administration est responsive et utilisable
  sur écrans desktop et tablette (résolution minimale~:
  1\,024~$\times$~768~px). & Moyenne \\
ENF10 & Extensibilité & Le schéma de données et l'API REST sont conçus pour
  accueillir de nouvelles régions synodales, types d'entités ou rôles sans
  refactorisation majeure du cœur du système. & Basse \\
\rowcolor{eeclightgreen!60}
ENF11 & Fiabilité & Le service de calcul d'itinéraire doit rester
  fonctionnel même en cas d'indisponibilité des services de routing
  externes~: une chaîne de repli à trois niveaux (OSRM, puis Valhalla, puis
  Haversine) garantit qu'un résultat exploitable est toujours renvoyé, sans
  jamais faire échouer la requête. & Haute \\
ENF12 & Performance temps réel & Pendant la navigation guidée~3D, la
  position de l'utilisateur doit être rafraîchie au moins une fois par
  seconde et le rendu cartographique doit rester fluide sur les
  navigateurs mobiles courants (WebGL matériel requis). & Moyenne \\
\end{longtable}

%------------------------------------------------------------
\subsection{Diagramme de Contexte}
\label{ssec:contexte}

Le diagramme de contexte (Figure~\ref{fig:diag-contexte}) positionne la
plateforme EEC~Géolocalisation dans son environnement en l'abordant comme une
boîte noire~: il identifie les acteurs externes, la nature de leurs interactions
avec le système et les flux d'information échangés, sans détailler les composants
internes. Six acteurs externes sont identifiés.

\begin{itemize}[itemsep=3pt, topsep=5pt, leftmargin=1.5cm]
  \item \textbf{Administrateur SUPER} (Bureau National)~: acteur le plus
    privilégié, il dispose d'un accès complet à toutes les données --- gestion
    de toutes les entités institutionnelles, import/export, consultation du
    journal d'audit, gestion des comptes utilisateurs et statistiques de
    toutes les régions synodales~;
  \item \textbf{Administrateur REGION}~: accède aux mêmes fonctions
    d'administration que le SUPER, mais dans les limites strictes de sa
    région synodale~; il peut également consulter les analytiques régionales~;
  \item \textbf{Administrateurs DISTRICT et PAROISSE}~: accèdent aux fonctions
    CRUD dans le périmètre de leur district ou paroisse respectif~;
  \item \textbf{Visiteur Non Authentifié}~: accède uniquement à la carte
    publique et aux informations générales des paroisses, sans aucun accès
    au module d'administration ni aux services personnalisés~;
  \item \textbf{Visiteur Authentifié}~: visiteur inscrit qui bénéficie de
    services supplémentaires depuis la carte publique --- enregistrement de
    paroisses en favoris, calcul d'itinéraire vers la paroisse la plus proche
    (PostGIS \texttt{Distance()}), et historique de ses recherches~;
  \item \textbf{GeoServer} (système cartographique)~: système externe qui reçoit
    les données géographiques depuis le backend Django et répond aux requêtes
    WMS~1.3.0 et WFS~2.0.0 émises par le frontend Leaflet.
\end{itemize}

\begin{figure}[H]
  \centering
  \hspace{1cm} % Déplace l'image vers la droite
  \includegraphics[width=1\textwidth]{diagrammes-Diagramme de Contexte.png}
  \caption{Diagramme de contexte du système de géolocalisation EEC}
  \label{fig:diag-contexte}
\end{figure}
%------------------------------------------------------------
\subsection{Diagramme de Packages}
\label{ssec:packages}

Le diagramme de packages (Figure~\ref{fig:diag-packages}) décompose le système
en grands modules et représente leurs dépendances. Trois couches principales
sont identifiées.

\noindent\textbf{Couche frontend (Next.js~16)~:}
\begin{itemize}[itemsep=2pt, topsep=3pt, leftmargin=1.5cm]
  \item \texttt{app/(public)}~: interface publique avec carte Leaflet, clustering,
    filtres et panneau de détail~;
  \item \texttt{app/admin/(general)}~: tableau de bord SUPER/REGION (paroisses,
    districts, régions, œuvres, ouvriers, statistiques, audit, comptes,
    import/export)~;
  \item \texttt{app/admin/district}~: tableau de bord DISTRICT (accès restreint
    au périmètre du district)~;
  \item \texttt{app/admin/paroisse}~: tableau de bord PAROISSE (accès limité à
    la paroisse de l'utilisateur).
\end{itemize}

\noindent\textbf{Couche backend (Django~5 + DRF)~:}
\begin{itemize}[itemsep=2pt, topsep=3pt, leftmargin=1.5cm]
  \item \texttt{apps/paroisses}, \texttt{apps/districts}, \texttt{apps/regions}~:
    modèles géospatiaux, sérialiseurs et endpoints CRUD avec filtrage RBAC~;
  \item \texttt{apps/oeuvres}, \texttt{apps/ouvriers}~: gestion des entités
    sociales et du personnel ecclésiastique~;
  \item \texttt{apps/statistiques}~: saisie et consultation des statistiques
    annuelles par paroisse~;
  \item \texttt{apps/comptes}~: gestion des comptes utilisateurs, authentification
    par sessions Django et middleware RBAC~;
  \item \texttt{apps/journal}~: enregistrement automatique via signaux Django
    de toutes les opérations d'écriture~;
  \item \texttt{apps/visiteurs}~: profil visiteur public, historique de
    recherche et itinéraires personnels.
\end{itemize}

\noindent\textbf{Couche infrastructure~:} PostgreSQL~16 + PostGIS~3.4 (persistance
spatiale), GeoServer~2.26 (publication WMS/WFS), Redis~7 (cache), Docker Compose
(orchestration des cinq services).

\begin{figure}[H]
  \center
  \includegraphics[width=1\textwidth]{diagrammes-Diagramme de package.png}
  \caption{Diagramme de contexte du système de géolocalisation EEC}
  \label{fig:diag-contexte}
\end{figure}

%------------------------------------------------------------
\subsection{Diagrammes de Cas d'Utilisation}
\label{ssec:uc}

Les diagrammes de cas d'utilisation modélisent les interactions entre les acteurs
et le système du point de vue fonctionnel. La Figure~\ref{fig:diag-uc} présente
le diagramme global consolidé de la plateforme, mettant en évidence les
vingt-neuf cas d'utilisation identifiés (UC01--UC29) et leur distribution
selon les six acteurs identifiés. Le tableau~\ref{tab:uc-recap} qui suit
détaille chacun de ces vingt-neuf cas d'utilisation --- intitulé, acteur
principal et description --- regroupés par périmètre d'accès~: public,
visiteur authentifié, administrateur paroissial, de district, régional,
national, et système cartographique.

\begin{figure}[H]
  \center
  \includegraphics[width=0.81\textwidth]{diagrammes-Diagramme de Cas d'Utilisation.11.png}
  \caption{Diagramme de contexte du système de géolocalisation EEC}
  \label{fig:diag-contexte}
\end{figure}

\begingroup
\small
\setlength{\tabcolsep}{5pt}
\renewcommand{\arraystretch}{1.20}

\begin{longtable}{%
  >{\bfseries\color{eccdarkgreen}}m{0.9cm}
  m{4.3cm}
  m{3.0cm}
  m{5.9cm}}

\caption{Récapitulatif des cas d'utilisation de la plateforme EEC}
\label{tab:uc-recap} \\

% ---- En-tête première page ----
\toprule
\rowcolor{eecgreen}
\multicolumn{1}{>{\bfseries\color{white}}m{0.9cm}}{ID} &
\multicolumn{1}{>{\bfseries\color{white}}m{4.3cm}}{Intitulé} &
\multicolumn{1}{>{\bfseries\color{white}}m{3.0cm}}{Acteur principal} &
\multicolumn{1}{>{\bfseries\color{white}}m{5.9cm}}{Description} \\
\midrule
\endfirsthead

% ---- En-tête pages suivantes ----
\multicolumn{4}{l}{\small\textit{(suite du Tableau~\ref{tab:uc-recap})}} \\[2pt]
\toprule
\rowcolor{eecgreen}
\multicolumn{1}{>{\bfseries\color{white}}m{0.9cm}}{ID} &
\multicolumn{1}{>{\bfseries\color{white}}m{4.3cm}}{Intitulé} &
\multicolumn{1}{>{\bfseries\color{white}}m{3.0cm}}{Acteur principal} &
\multicolumn{1}{>{\bfseries\color{white}}m{5.9cm}}{Description} \\
\midrule
\endhead

% ---- Pied de page intermédiaire ----
\midrule
\multicolumn{4}{r}{\small\textit{(suite page suivante\ldots)}} \\
\endfoot

% ---- Pied de dernière page ----
\bottomrule
\endlastfoot

% ==================================================
\multicolumn{4}{>{\bfseries\color{white}}m{\linewidth}}%
  {\cellcolor{eccdarkgreen}\quad Accès public --- aucune connexion requise} \\

UC01 & Consulter la carte interactive
  & Tout utilisateur
  & Afficher la carte générale avec les marqueurs des paroisses et des
    œuvres de l'EEC, regroupés automatiquement selon le niveau de zoom. \\
\rowcolor{eeclightgreen!60}
UC02 & Rechercher une entité
  & Tout utilisateur
  & Trouver une paroisse ou une œuvre par son nom depuis la barre de
    recherche de la carte. \\
UC03 & Afficher la fiche d'une entité
  & Tout utilisateur
  & Consulter les informations détaillées d'une paroisse ou d'une œuvre
    en sélectionnant son marqueur sur la carte. \\
\rowcolor{eeclightgreen!60}
UC04 & Localiser la paroisse la plus proche
  & Tout utilisateur
  & Identifier la paroisse géographiquement la plus proche d'une
    position donnée. \\
UC05 & S'authentifier
  & Tout utilisateur enregistré
  & Accéder à la plateforme en saisissant ses identifiants et ouvrir
    une session personnelle sécurisée. \\

% ==================================================
\multicolumn{4}{>{\bfseries\color{white}}m{\linewidth}}%
  {\cellcolor{eccdarkgreen}\quad Visiteur authentifié} \\

\rowcolor{eeclightgreen!60}
UC06 & Gérer son profil
  & Visiteur authentifié
  & Consulter et mettre à jour ses informations personnelles et ses
    préférences de navigation. \\
UC07 & Consulter son historique
  & Visiteur authentifié
  & Accéder à la liste des dernières paroisses consultées et des
    recherches précédemment effectuées. \\
\rowcolor{eeclightgreen!60}
UC08 & Effacer son historique
  & Visiteur authentifié
  & Supprimer l'ensemble des entrées enregistrées dans son historique
    de navigation. \\
UC09 & Enregistrer un favori
  & Visiteur authentifié
  & Sauvegarder une paroisse dans sa liste personnelle pour un accès
    rapide ultérieur. \\
\rowcolor{eeclightgreen!60}
UC10 & Retirer un favori
  & Visiteur authentifié
  & Supprimer une paroisse précédemment enregistrée de sa liste de
    favoris. \\
UC11 & Changer son mot de passe
  & Visiteur authentifié
  & Modifier son mot de passe de connexion, obligatoirement demandé
    lors de la toute première connexion au compte. \\
\rowcolor{eeclightgreen!60}
UC12 & Calculer un itinéraire
  & Visiteur authentifié
  & Obtenir la distance et la durée de trajet estimée entre un point
    de départ et une paroisse de destination choisie. \\

% ==================================================
\multicolumn{4}{>{\bfseries\color{white}}m{\linewidth}}%
  {\cellcolor{eccdarkgreen}\quad Administrateur paroissial} \\

UC13 & Gérer la fiche de sa paroisse
  & Admin paroissial
  & Modifier les informations générales et les coordonnées
    géographiques de la paroisse dont il est responsable. \\
\rowcolor{eeclightgreen!60}
UC14 & Gérer les ouvriers
  & Admin paroissial
  & Créer, modifier et affecter les membres du personnel
    ecclésiastique rattachés à sa paroisse. \\
UC15 & Gérer les œuvres
  & Admin paroissial
  & Administrer les établissements sociaux et éducatifs rattachés
    à sa paroisse. \\
\rowcolor{eeclightgreen!60}
UC16 & Saisir les statistiques annuelles
  & Admin paroissial
  & Enregistrer les données chiffrées de l'année pour la paroisse~:
    effectifs, cérémonies et indicateurs financiers. \\
UC17 & Exporter les données
  & Admin paroissial
  & Télécharger la liste des entités de sa paroisse au format tableur
    ou au format document imprimable. \\
\rowcolor{eeclightgreen!60}
UC18 & Consulter le tableau de bord
  & Admin paroissial
  & Visualiser les indicateurs de synthèse de sa paroisse~: effectifs,
    taux de couverture et évolution statistique par année. \\
UC19 & Gérer son compte
  & Admin paroissial
  & Consulter et modifier ses informations de connexion et les
    paramètres de son compte administrateur. \\
\rowcolor{eeclightgreen!60}
UC20 & Réinitialiser le mot de passe
  & Admin paroissial
  & Modifier le mot de passe attribué par défaut lors de la première
    connexion à la plateforme. \\

% ==================================================
\multicolumn{4}{>{\bfseries\color{white}}m{\linewidth}}%
  {\cellcolor{eccdarkgreen}\quad Administrateur de district} \\

UC21 & Gérer les paroisses du district
  & Admin district
  & Créer, modifier et supprimer les paroisses situées dans le
    périmètre de son district ecclésial. \\
\rowcolor{eeclightgreen!60}
UC22 & Valider les données saisies
  & Admin district
  & Vérifier la cohérence hiérarchique et géographique des
    informations avant leur enregistrement définitif. \\

% ==================================================
\multicolumn{4}{>{\bfseries\color{white}}m{\linewidth}}%
  {\cellcolor{eccdarkgreen}\quad Administrateur régional} \\

UC23 & Gérer les entités de la région
  & Admin régional
  & Administrer l'ensemble des districts, paroisses et entités
    rattachés à sa région synodale. \\

% ==================================================
\multicolumn{4}{>{\bfseries\color{white}}m{\linewidth}}%
  {\cellcolor{eccdarkgreen}\quad Administrateur national} \\

\rowcolor{eeclightgreen!60}
UC24 & Administrer toutes les entités
  & Admin national
  & Créer, modifier et supprimer toute entité de l'EEC sur l'ensemble
    du territoire, sans restriction géographique. \\
UC25 & Gérer les comptes utilisateurs
  & Admin national
  & Créer, activer, désactiver et paramétrer les comptes de tous
    les administrateurs de la plateforme. \\
\rowcolor{eeclightgreen!60}
UC26 & Importer les données nationales
  & Admin national
  & Charger en masse les données institutionnelles depuis les fichiers
    sources fournis par les services de l'EEC. \\
UC27 & Consulter le journal d'audit
  & Admin national
  & Accéder à l'historique complet de toutes les modifications
    effectuées sur les données de la plateforme. \\

% ==================================================
\multicolumn{4}{>{\bfseries\color{white}}m{\linewidth}}%
  {\cellcolor{eccdarkgreen}\quad Système cartographique} \\

UC28 & Publier les couches cartographiques
  & Système cartographique
  & Diffuser les données géographiques de l'EEC sous forme de couches
    superposables accessibles depuis la carte interactive. \\
\rowcolor{eeclightgreen!60}
UC29 & Répondre aux requêtes de la carte
  & Système cartographique
  & Traiter les demandes d'images et de données géographiques émises
    par la carte lors de la navigation de l'utilisateur. \\

\end{longtable}
\endgroup

%------------------------------------------------------------
\subsection{Description Textuelle des Cas d'Utilisation Principaux}
\label{ssec:uc-detail}

Parmi les vingt-neuf cas d'utilisation recensés au tableau~\ref{tab:uc-recap},
dix ont été retenus comme représentatifs de l'ensemble des profils et des
domaines fonctionnels de la plateforme~: consultation cartographique,
recherche et localisation, authentification, services du visiteur
authentifié, administration des entités, export et pilotage. Chacun est
décrit ci-après selon un format structuré~: acteur principal, préconditions,
postconditions et scénario nominal.

\begin{ucfiche}{UC01}{Consulter la carte interactive}
\textbf{Acteur~:} Tout utilisateur (y compris non authentifié).\\[2pt]
\textbf{Préconditions~:} Aucune~; la carte publique est accessible sans
connexion.\\[2pt]
\textbf{Postconditions~:} La carte affiche les marqueurs des paroisses et
des œuvres regroupés en clusters selon le niveau de zoom courant.\\[2pt]
\textbf{Scénario nominal~:}
\begin{enumerate}[itemsep=1pt, topsep=2pt, leftmargin=1.2cm]
  \item L'utilisateur accède à la page d'accueil publique de la plateforme.
  \item Le frontend interroge GeoServer (WMS) et l'API backend (GeoJSON) pour
    récupérer les entités géolocalisées.
  \item Leaflet affiche la carte avec les couches institutionnelles
    superposées au fond OpenStreetMap, marqueurs regroupés en clusters.
  \item L'utilisateur navigue (zoom, déplacement)~; les clusters se
    subdivisent dynamiquement en marqueurs individuels.
\end{enumerate}
\end{ucfiche}

\begin{ucfiche}{UC02}{Rechercher une entité}
\textbf{Acteur~:} Tout utilisateur.\\[2pt]
\textbf{Préconditions~:} La carte est chargée (UC01).\\[2pt]
\textbf{Postconditions~:} La carte se recentre sur l'entité trouvée et
affiche son marqueur mis en évidence.\\[2pt]
\textbf{Scénario nominal~:}
\begin{enumerate}[itemsep=1pt, topsep=2pt, leftmargin=1.2cm]
  \item L'utilisateur saisit un nom (paroisse, œuvre) dans la barre de
    recherche.
  \item Le frontend interroge l'API avec le terme saisi, combiné le cas
    échéant aux filtres actifs (région, district, classe de paroisse).
  \item L'API retourne la liste des entités correspondantes, triée par
    pertinence.
  \item L'utilisateur sélectionne un résultat~; la carte se recentre et
    ouvre la popup de détail correspondante.
\end{enumerate}
\end{ucfiche}

\begin{ucfiche}{UC03}{Localiser la paroisse la plus proche}
\textbf{Acteur~:} Tout utilisateur.\\[2pt]
\textbf{Préconditions~:} L'utilisateur autorise la géolocalisation de son
navigateur, ou clique sur un point de la carte.\\[2pt]
\textbf{Postconditions~:} La paroisse géographiquement la plus proche de la
position de référence est identifiée et mise en évidence.\\[2pt]
\textbf{Scénario nominal~:}
\begin{enumerate}[itemsep=1pt, topsep=2pt, leftmargin=1.2cm]
  \item L'utilisateur déclenche la géolocalisation ou sélectionne un point
    de départ sur la carte.
  \item Le backend calcule, via la fonction spatiale
    \texttt{ST\_Distance()} de PostGIS, la paroisse géolocalisée la plus
    proche de ce point.
  \item Le résultat (nom, distance à vol d'oiseau) est renvoyé et affiché à
    l'utilisateur.
\end{enumerate}
\end{ucfiche}

\begin{ucfiche}{UC04}{S'authentifier}
\textbf{Acteur~:} Tout utilisateur enregistré (administrateur ou visiteur).\\[2pt]
\textbf{Préconditions~:} L'utilisateur dispose d'un compte actif.\\[2pt]
\textbf{Postconditions~:} Une session sécurisée est ouverte~; l'utilisateur
est redirigé vers l'interface correspondant à son rôle.\\[2pt]
\textbf{Scénario nominal~:}
\begin{enumerate}[itemsep=1pt, topsep=2pt, leftmargin=1.2cm]
  \item L'utilisateur accède au formulaire de connexion et saisit ses
    identifiants.
  \item Le frontend récupère un jeton CSRF puis soumet les identifiants à
    l'API d'authentification.
  \item Le backend vérifie les identifiants et ouvre une session Django
    (cookie \textit{HttpOnly}, \textit{Secure}, \textit{SameSite}).
  \item L'utilisateur est redirigé vers le tableau de bord de son rôle~;
    si la première connexion est détectée, il est redirigé vers le
    changement de mot de passe obligatoire.
\end{enumerate}
\textbf{Scénario alternatif~:} identifiants invalides~$\rightarrow$ message
d'erreur générique, sans révéler si l'email ou le mot de passe est erroné
(bonne pratique OWASP~[11]).
\end{ucfiche}

\begin{ucfiche}{UC05}{Consulter l'historique des recherches}
\textbf{Acteur~:} Visiteur authentifié.\\[2pt]
\textbf{Préconditions~:} L'utilisateur est connecté (UC04) et a effectué au
moins une recherche ou un calcul d'itinéraire antérieur.\\[2pt]
\textbf{Postconditions~:} La liste chronologique des recherches passées est
affichée dans l'espace personnel.\\[2pt]
\textbf{Scénario nominal~:}
\begin{enumerate}[itemsep=1pt, topsep=2pt, leftmargin=1.2cm]
  \item Le visiteur accède à son profil personnel, onglet \og Historique \fg{}.
  \item Le frontend interroge l'API de l'historique du visiteur connecté.
  \item Les entrées sont affichées par ordre chronologique inversé (les plus
    récentes en premier).
  \item Le visiteur peut effacer tout ou partie de son historique.
\end{enumerate}
\end{ucfiche}

\begin{ucfiche}{UC06}{Créer une paroisse}
\textbf{Acteur~:} Administrateur SUPER, REGION ou DISTRICT (selon son
périmètre synodal).\\[2pt]
\textbf{Préconditions~:} L'administrateur est authentifié (UC04) et dispose
des droits de création dans le périmètre concerné.\\[2pt]
\textbf{Postconditions~:} La nouvelle paroisse est enregistrée en base,
géolocalisée si les coordonnées sont fournies, et une entrée d'audit
CREATE est journalisée.\\[2pt]
\textbf{Scénario nominal~:}
\begin{enumerate}[itemsep=1pt, topsep=2pt, leftmargin=1.2cm]
  \item L'administrateur ouvre le formulaire de création de paroisse et
    saisit les informations générales (nom, district, classe/catégorie).
  \item Il positionne la localisation GPS via la carte interactive
    intégrée au formulaire.
  \item Le frontend valide les champs obligatoires et soumet la requête
    \texttt{POST} à l'API.
  \item Le backend applique le contrôle RBAC (le district doit appartenir
    au périmètre de l'administrateur), valide les coordonnées et
    enregistre la paroisse.
  \item Une entrée d'audit est créée automatiquement (signal Django).
\end{enumerate}
\end{ucfiche}

\begin{ucfiche}{UC07}{Exporter les données}
\textbf{Acteur~:} Administrateur SUPER, REGION, DISTRICT ou PAROISSE.\\[2pt]
\textbf{Préconditions~:} L'administrateur est authentifié et consulte une
liste d'entités (paroisses, œuvres, ouvriers).\\[2pt]
\textbf{Postconditions~:} Un fichier Excel ou PDF contenant les données du
périmètre de l'utilisateur est généré et téléchargé.\\[2pt]
\textbf{Scénario nominal~:}
\begin{enumerate}[itemsep=1pt, topsep=2pt, leftmargin=1.2cm]
  \item L'administrateur sélectionne le format d'export souhaité (Excel ou
    PDF) depuis l'interface de liste.
  \item Le backend vérifie le rôle de l'utilisateur puis extrait les
    données de son périmètre par une requête SQL agrégée unique.
  \item Le fichier est généré (\texttt{openpyxl} pour Excel,
    \texttt{WeasyPrint} pour PDF) selon la charte graphique EEC.
  \item Le fichier est renvoyé au navigateur et l'opération est journalisée
    dans le journal d'audit.
\end{enumerate}
\end{ucfiche}

\begin{ucfiche}{UC08}{Consulter le tableau de bord}
\textbf{Acteur~:} Tout administrateur (SUPER, REGION, DISTRICT, PAROISSE).\\[2pt]
\textbf{Préconditions~:} L'administrateur est authentifié (UC04).\\[2pt]
\textbf{Postconditions~:} Les indicateurs de synthèse du périmètre de
l'utilisateur sont affichés.\\[2pt]
\textbf{Scénario nominal~:}
\begin{enumerate}[itemsep=1pt, topsep=2pt, leftmargin=1.2cm]
  \item L'administrateur accède à la page d'accueil de son espace
    d'administration.
  \item Le backend calcule dynamiquement les indicateurs (nombre de
    paroisses, districts, ouvriers, œuvres, taux de couverture GPS) filtrés
    selon le rôle et le périmètre de l'utilisateur.
  \item Le tableau de bord affiche ces indicateurs, ainsi que l'évolution
    statistique sur les dernières années.
\end{enumerate}
\end{ucfiche}

\begin{ucfiche}{UC09}{Calculer un itinéraire}
\textbf{Acteur~:} Visiteur authentifié.\\[2pt]
\textbf{Préconditions~:} L'utilisateur est connecté (UC04) et a défini un
point de départ (sa position GPS ou une entité) et une destination.\\[2pt]
\textbf{Postconditions~:} Un itinéraire réel (distance, durée, tracé,
étapes) est calculé et affiché sur la carte~; le visiteur peut ensuite
démarrer la navigation guidée en~3D.\\[2pt]
\textbf{Scénario nominal~:}
\begin{enumerate}[itemsep=1pt, topsep=2pt, leftmargin=1.2cm]
  \item Le visiteur ouvre l'onglet \og Parcours \fg{}, définit départ,
    destination et mode de déplacement (voiture, vélo, à pied).
  \item Le frontend envoie une requête \texttt{POST~/api/visitor/itineraires/route/}.
  \item Le backend interroge successivement OSRM puis Valhalla
    (section~\ref{ssec:justification-algo})~; en cas de double échec, un
    calcul Haversine est appliqué en repli.
  \item Le tracé, la distance et la durée sont renvoyés et dessinés sur la
    carte Leaflet, avec les instructions virage-par-virage.
\end{enumerate}
\textbf{Scénario alternatif~:} le visiteur clique sur \og Démarrer la
navigation \fg{}~$\rightarrow$ enchaîne vers la vue~3D guidée (suivi GPS en
temps réel, décrite en section~\ref{ssec:res-avance}).
\end{ucfiche}

\begin{ucfiche}{UC10}{Gérer les comptes utilisateurs}
\textbf{Acteur~:} Administrateur SUPER (tout compte)~; administrateurs
REGION et DISTRICT (comptes de leur périmètre uniquement).\\[2pt]
\textbf{Préconditions~:} L'administrateur est authentifié et dispose des
droits de gestion des comptes.\\[2pt]
\textbf{Postconditions~:} Le compte administrateur est créé, modifié ou
désactivé, avec un mot de passe temporaire à changer à la première
connexion.\\[2pt]
\textbf{Scénario nominal~:}
\begin{enumerate}[itemsep=1pt, topsep=2pt, leftmargin=1.2cm]
  \item L'administrateur accède à la gestion des comptes et sélectionne
    \og Créer un administrateur \fg{}.
  \item Il saisit les informations du compte (identité, rôle, périmètre
    synodal --- limité à son propre périmètre pour REGION et DISTRICT).
  \item Le backend applique le contrôle RBAC hiérarchique~: un
    administrateur REGION ne peut créer que des comptes DISTRICT ou
    PAROISSE de sa propre région.
  \item Le compte est créé avec un mot de passe temporaire~; l'indicateur
    de première connexion force son changement dès la prochaine connexion.
\end{enumerate}
\end{ucfiche}

%------------------------------------------------------------
\subsection{Diagramme de Classes Métier}
\label{ssec:classes-metier}

Le diagramme de classes métier    modélise
les entités du domaine institutionnel de l'EEC et leurs relations, indépendamment
de tout détail d'implémentation technique. Huit entités principales sont
identifiées~:


\begin{figure}[H]
  \center
  \includegraphics[width=0.81\textwidth]{diagrammes- classe metier.png}
  \caption{Diagramme de classes métier du domaine institutionnel EEC}
  \label{fig:diag-contexte}
\end{figure}

%------------------------------------------------------------
\subsection{Diagramme d'Activité}
\label{ssec:activite}

Le diagramme d'activité modélise le flux de traitement du processus d'import
de données depuis Excel tel que conçu et implémenté lors de la phase initiale
du projet --- le processus métier le plus critique et le plus complexe du
système à cette étape, car il constituait alors le seul point d'entrée de
masse pour les données institutionnelles et impliquait une chaîne de
validation multi-étapes avant toute persistance. L'audit approfondi des
données sources (section~\ref{sec:audit-donnees}) a toutefois révélé des
incohérences trop nombreuses pour être fiabilisées par une validation
automatique~: ce module a donc été retiré de la version finale au profit
d'une saisie directement contrôlée dans l'interface d'administration. Ce
diagramme est conservé ici à titre de documentation de conception, la
logique de validation multi-étapes qu'il décrit ayant directement inspiré
les contrôles de saisie de l'interface actuelle.

Le flux se déroule en cinq phases~: (1)~téléversement et validation du format
du fichier Excel~; (2)~extraction et validation ligne par ligne (coordonnées GPS,
cohérence hiérarchique, doublons)~; (3)~insertion des lignes valides en base de
données PostGIS~; (4)~enregistrement d'une entrée d'audit synthétique~;
(5)~génération et affichage du rapport d'import détaillé.

\begin{figure}[H]
  \center
  \includegraphics[width=0.55\textwidth]{importerdonner.png}
  \caption{Diagramme d’activité du processus d’import de données depuis Exce}
  \label{fig:diag-contexte}
\end{figure}



%============================================================
\section{Conception du Système}
\label{sec:conception}
%============================================================

\subsection{Architecture Technique}
\label{ssec:archi}

La Figure si dessous présente l'architecture technique retenue~: une
architecture à cinq services découplés, orchestrés par Docker Compose sur un
serveur unique. Chaque service est isolé dans son propre conteneur et communique
avec les autres via le réseau Docker interne \texttt{eec\_network}.

\noindent\textbf{Services et rôles~:}
\begin{itemize}[itemsep=2pt, topsep=4pt, leftmargin=1.5cm]
  \item \textbf{Frontend Next.js~16} (port~3004)~: application React côté
    serveur avec App Router, TypeScript et Tailwind CSS. Sert la carte publique
    Leaflet et les quatre interfaces d'administration. Communique avec le backend
    exclusivement via l'API REST (\texttt{/api/}), sans accès direct à la base
    de données~;
  \item \textbf{Backend Django~5 + DRF} (port~8000)~: API REST de référence.
    Expose 40+ endpoints (geo, auth, exports, imports, visiteurs, analytics,
    audit). Applique le RBAC à chaque requête via un middleware dédié. Gère les
    sessions dans Redis. Exécute les requêtes géospatiales via GeoDjango et
    PostGIS~;
  \item \textbf{PostgreSQL~16 + PostGIS~3.4} (port~5432)~: base de données
    relationnelle et spatiale. Stocke l'intégralité des données institutionnelles
    (paroisses avec \texttt{PointField}, régions avec \texttt{MultiPolygonField})
    et les données métier (ouvriers, œuvres, statistiques, audit)~;
  \item \textbf{GeoServer~2.26} (port~8080)~: serveur cartographique OGC. Publie
    les couches géographiques de la base PostGIS en WMS~1.3.0 (images raster)
    et WFS~2.0.0 (features vectorielles). Consommé directement par Leaflet dans
    le frontend~;
  \item \textbf{Redis~7} (port~6379)~: stockage en mémoire utilisé pour la
    gestion des sessions Django (\texttt{SESSION\_ENGINE =
    django.contrib.sessions.backends.cache}) et la mise en cache des requêtes
    fréquentes.
\end{itemize}

\noindent\textbf{Flux de communication~:} Le frontend Next.js adresse ses
requêtes API au backend Django. Le backend interroge PostgreSQL/PostGIS pour
toutes les opérations de données et Redis pour les sessions. GeoServer accède
directement à PostGIS via JDBC pour lire les couches géographiques et les
servir au frontend Leaflet. Aucun service n'est exposé directement à Internet
--- un reverse proxy (Nginx) centralise les accès externes en production.

\begin{figure}[H]
  \center
  \includegraphics[width=0.81\textwidth]{archtecture technique.png}
  \caption{Architecture technique de la plateforme EEC Géolocalisation}
  \label{fig:diag-contexte}
\end{figure}

%------------------------------------------------------------
\subsection{Modèle de Données}
\label{ssec:modele-bd}

Le modèle de données repose sur quinze tables PostGIS réparties en cinq
domaines fonctionnels. Le Tableau~\ref{tab:modele-bd} synthétise les tables
principales, leurs champs géospatiaux et leurs relations clés~; la
Figure~\ref{fig:diag-modele-bd} en présente le schéma entité-relation.

\setlength{\tabcolsep}{5pt}
\renewcommand{\arraystretch}{1.22}
\begin{longtable}{%
  >{\bfseries\color{eccdarkgreen}}m{3.5cm}
  m{4.5cm}
  m{6.3cm}}
\caption{Principales tables du modèle de données PostGIS}\label{tab:modele-bd} \\
\toprule
\rowcolor{eecgreen}
\multicolumn{1}{>{\bfseries\color{white}}m{3.5cm}}{Table} &
\multicolumn{1}{>{\bfseries\color{white}}m{4.5cm}}{Champ géospatial / clé} &
\multicolumn{1}{>{\bfseries\color{white}}m{6.3cm}}{Rôle et relations principales} \\
\midrule
\endfirsthead
\caption[]{\textit{(suite)~Modèle de données PostGIS}} \\
\toprule
\rowcolor{eecgreen}
\multicolumn{1}{>{\bfseries\color{white}}m{3.5cm}}{Table} &
\multicolumn{1}{>{\bfseries\color{white}}m{4.5cm}}{Champ géospatial / clé} &
\multicolumn{1}{>{\bfseries\color{white}}m{6.3cm}}{Rôle et relations principales} \\
\midrule
\endhead
\midrule
\multicolumn{3}{r}{\small\textit{(suite page suivante)}} \\
\endfoot
\bottomrule
\endlastfoot
\rowcolor{eeclightgreen!60}
geo\_regionsynodale & \texttt{polygone} MultiPolygonField (SRID=4326) & Entité de niveau 1~; racine de la hiérarchie synodale \\
geo\_district & FK région (PROTECT) & Entité de niveau 2~; 159 districts \\
\rowcolor{eeclightgreen!60}
geo\_paroisse & \texttt{position} PointField (SRID=4326) & Entité de base~; FK district (PROTECT)~; \texttt{categorie} (classe~: C1--C4, B1--B2, A1--A2, A++)~; 858 paroisses \\
oeuvres\_oeuvre & \texttt{localisation} PointField (SRID=4326) & Établissements sociaux/éducatifs~; FK paroisse (SET\_NULL) \\
\rowcolor{eeclightgreen!60}
ouvriers\_ouvrier & --- & Personnel ecclésiastique~; FK paroisse (SET\_NULL) \\
statistiques\_stat & --- & Relevé annuel par paroisse~; FK paroisse (CASCADE) \\
\rowcolor{eeclightgreen!60}
accounts\_utilisateur & --- & Étend AbstractUser~; champs role, region, district, paroisse \\
accounts\_profil & --- & OneToOne vers Utilisateur~; informations de contact étendues \\
\rowcolor{eeclightgreen!60}
audit\_entreeaudit & --- & Journal immuable~; FK utilisateur (SET\_NULL) \\
visitors\_profilvisiteur & --- & OneToOne vers User~; créé auto à l'inscription \\
\rowcolor{eeclightgreen!60}
visitors\_paroisseenreg & --- & M2M Visiteur × Paroisse~; date d'enregistrement \\
visitors\_itineraire & \texttt{point\_depart} PointField & FK ProfilVisiteur (CASCADE)~; paroisse cible FK \\
\rowcolor{eeclightgreen!60}
visitors\_recherche & --- & Historique de recherche~; FK ProfilVisiteur (CASCADE) \\
exports\_rapport & --- & Métadonnées des exports générés (PDF, XLSX) \\
\rowcolor{eeclightgreen!60}
geo\_paroissevue & --- & Compteur de vues par paroisse~; FK paroisse (CASCADE) \\
\end{longtable}

\subsection{Diagramme de Séquence}
\label{ssec:sequence}

Le diagramme de séquence    modélise le flux
d'authentification sécurisée, scénario critique qui conditionne l'accès à
l'ensemble des fonctionnalités protégées. Ce flux illustre les interactions
entre le navigateur, le frontend Next.js, le middleware Django et la base
de données Redis.

Le diagramme de séquence technique cas  d untlisation :consulter la carte

détaille les échanges
techniques entre les composants lors de la consultation de la carte interactive.
Lors du chargement initial, le navigateur interroge successivement GeoServer
pour obtenir les capacités WMS et les tuiles cartographiques, puis l'API Django
pour récupérer les paroisses géolocalisées au format GeoJSON, qui sont ensuite
regroupées automatiquement en clusters par Leaflet.

Lorsque l'utilisateur applique un filtre hiérarchique, une requête filtrée est
envoyée à l'API (\texttt{GET /api/geo/paroisses/?region=\{id\}}), Django
exécute la requête SQL correspondante sur PostGIS et retourne un GeoJSON mis
à jour. Un clic sur un marqueur déclenche l'instanciation d'un objet
\textit{popup} dans l'interface Leaflet, affichant les informations complètes
de la paroisse sans requête supplémentaire (données déjà chargées dans le
GeoJSON initial).


 


\begin{figure}[H]
  \center
  \includegraphics[width=1\textwidth]{technique1.png}
  \caption{ digaramme sequence technique :consulter la carte}
  \label{fig:diag-contexte}
\end{figure}

\subsection{Diagramme de sequence technique cas d utilisation :exporterdocument   }
\label{ssec:classes-tech}
%------------------------------------------------------------
Les diagrammes de séquence des Figures~\ref{fig:seq-export-excel}
et~\ref{fig:seq-export-pdf} modélisent le flux technique d'export de
données. Dans les deux cas, le middleware RBAC intercepte la requête
avant tout accès aux données~: seuls les rôles habilités
(SUPER, REGION, DISTRICT, PAROISSE) peuvent déclencher un export.
Les données sont extraites de PostGIS en une seule requête SQL
agrégée, puis transmises au générateur de fichier. Pour le format
Excel, \texttt{openpyxl} applique la charte graphique EEC
(en-tête vert foncé, lignes alternées, colonnes gelées). Pour le
format PDF, \texttt{WeasyPrint} compose le document en A4 paysage
depuis un template HTML. Toute opération d'export est systématiquement
enregistrée dans le journal d'audit via \texttt{log\_action()}.


\begin{figure}[H]
  \center
  \includegraphics[width=1.1\textwidth]{technique2.png}
  \caption{ digaramme sequence technique :exporter pdf}
  \label{fig:diag-contexte}
\end{figure}

%------------------------------------------------------------
\subsection{Diagramme de Classes Technique}
\label{ssec:classes-tech}

Le diagramme de classes technique  
représente l'implémentation concrète des entités sous forme de modèles Django,
avec leurs attributs, types Python/PostGIS et relations inter-classes.

\noindent\textbf{Modèles principaux et attributs~:}
\begin{itemize}[itemsep=2pt, topsep=4pt, leftmargin=1.5cm]
  \item \textbf{Paroisse} (app \texttt{geo})~: \texttt{nom} (CharField),
    \texttt{district} (FK → District, PROTECT), \texttt{localisation}
    (PointField, SRID=4326, null=True), \texttt{statut} (CharField, choix
    ACTIVE/INACTIVE), \texttt{has\_gps} (BooleanField)~;
  \item \textbf{RegionSynodale} (app \texttt{geo})~: \texttt{nom} (CharField),
    \texttt{siege} (CharField), \texttt{polygone} (MultiPolygonField,
    SRID=4326)~;
  \item \textbf{Utilisateur} (app \texttt{accounts})~: étend
    \texttt{AbstractUser}~; \texttt{role} (CharField, choix SUPER/REGION/
    DISTRICT/PAROISSE/VISITEUR), \texttt{region} (FK → RegionSynodale,
    SET\_NULL), \texttt{district} (FK → District, SET\_NULL), \texttt{paroisse}
    (FK → Paroisse, SET\_NULL)~;
  \item \textbf{ProfilVisiteur} (app \texttt{visitors})~: \texttt{user}
    (OneToOneField → Utilisateur, CASCADE), \texttt{date\_inscription}
    (DateTimeField, auto), \texttt{derniere\_position} (PointField, null=True)~;
  \item \textbf{ItinerairePersonnel} (app \texttt{visitors})~: \texttt{profil}
    (FK → ProfilVisiteur, CASCADE), \texttt{paroisse\_cible} (FK → Paroisse,
    CASCADE), \texttt{point\_depart} (PointField), \texttt{distance\_km}
    (FloatField), \texttt{duree\_min} (IntegerField)~;
  \item \textbf{EntreeAudit} (app \texttt{audit})~: \texttt{auteur} (FK →
    Utilisateur, SET\_NULL), \texttt{action} (CharField, choix CREATE/UPDATE/
    DELETE), \texttt{modele} (CharField), \texttt{objet\_id} (IntegerField),
    \texttt{horodatage} (DateTimeField, auto\_now\_add).
\end{itemize}

\noindent\textbf{Relations inter-classes~:} la relation Paroisse → District →
RegionSynodale forme la colonne vertébrale hiérarchique (PROTECT empêche la
suppression en cascade). Les relations Oeuvre et Ouvrier vers Paroisse sont
SET\_NULL (un ouvrier peut exister sans paroisse). EntreeAudit est reliée à
Utilisateur en SET\_NULL pour conserver les logs même après suppression du
compte. ProfilVisiteur est en OneToOne avec Utilisateur (CASCADE).

\begin{figure}[H]
  \center
  \includegraphics[width=1.15\textwidth]{diagrammes-classetechnique.png}
  \caption{Diagramme de classes technique  }
  \label{fig:diag-classes-tech}
\end{figure}
%-------------

%============================================================
\section{Modélisation Mathématique et Algorithmique du Calcul d'Itinéraire}
\label{sec:modelisation-math}
%============================================================

Au-delà de la gestion des entités institutionnelles, la plateforme offre au
visiteur authentifié un service de calcul d'itinéraire réel entre deux points
géographiques (paroisse, œuvre, ou position GPS de l'utilisateur). Ce service ne
se limite pas à une mesure de distance~: il repose sur une modélisation
mathématique rigoureuse du réseau routier sous forme de graphe et sur des
algorithmes de recherche de plus court chemin. Cette section formalise cette
modélisation, indépendamment de son implémentation technique détaillée en
section~\ref{ssec:justification-algo} du Chapitre~3.

\subsection{Le Réseau Routier comme Graphe Pondéré}
\label{ssec:modele-graphe}

\subsubsection{Définition formelle}

Le réseau routier est modélisé comme un graphe orienté pondéré
$G = (V, E, w)$, où~:
\begin{itemize}[itemsep=2pt, topsep=4pt, leftmargin=1.5cm]
  \item $V$ est l'ensemble des \textbf{sommets} (\textit{vertices})~: chaque
    sommet $v \in V$ représente une intersection ou un carrefour du réseau
    routier camerounais~;
  \item $E \subseteq V \times V$ est l'ensemble des \textbf{arêtes}
    (\textit{edges})~: chaque arête $(u, v) \in E$ représente un tronçon de
    route reliant directement deux intersections, praticable dans le sens
    $u \rightarrow v$~;
  \item $w : E \rightarrow \mathbb{R}^{+}$ est la \textbf{fonction de poids}~:
    à chaque arête $(u, v)$ est associé un coût $w(u, v) > 0$, exprimé en
    temps de parcours plutôt qu'en simple distance.
\end{itemize}

Ce graphe est \textbf{orienté} (une rue à sens unique n'autorise le
déplacement que dans une direction) et ses poids sont \textbf{strictement
positifs}, propriété indispensable à la validité des algorithmes présentés en
section~\ref{ssec:dijkstra-astar}. Dans notre application, ce graphe n'est
pas construit ni maintenu par nos soins~: il est fourni par les données
\textbf{OpenStreetMap} et exploité par le moteur de routing externe (voir
section~\ref{ssec:justification-algo}).

\subsubsection{Fonction de poids~: le coût temporel}

Le poids d'une arête $(u, v)$ n'est pas sa seule longueur $d(u, v)$~: c'est le
\textbf{temps de parcours} de ce tronçon, qui dépend du mode de déplacement
choisi par l'utilisateur~: voiture, vélo ou piéton. Formellement~:

\begin{equation}
  w(u, v) = \frac{d(u, v)}{v_{\text{max}}(u, v, m)}
  \label{eq:poids-arete}
\end{equation}

où $d(u, v)$ est la longueur du tronçon (en kilomètres), et
$v_{\text{max}}(u, v, m)$ est la vitesse maximale autorisée sur ce tronçon
pour le mode $m \in \{\text{auto}, \text{bicycle}, \text{pedestrian}\}$. Cette
formulation a une conséquence importante~: \textbf{le graphe effectif change
avec le mode de transport}. Une arête correspondant à une autoroute a un
poids faible en voiture mais un poids infini ($w \rightarrow \infty$) pour le
mode piéton, puisque l'autoroute y est interdite~; elle disparaît donc
effectivement du sous-graphe accessible au piéton.

\subsubsection{Formulation du problème}

Étant donné un sommet de départ $s \in V$ (position de l'utilisateur ou
entité choisie) et un sommet d'arrivée $t \in V$ (entité de destination), le
problème à résoudre est celui du \textbf{plus court chemin}~: trouver la
séquence de sommets $s = v_0, v_1, \dots, v_k = t$ telle que la somme des
poids des arêtes traversées soit minimale~:

\begin{equation}
  \min_{\substack{(v_0, \dots, v_k) \\ \text{chemin de } s \text{ à } t}}
  \sum_{i=0}^{k-1} w(v_i, v_{i+1})
  \label{eq:plus-court-chemin}
\end{equation}

Ce problème d'optimisation combinatoire porte le nom de \textit{Single-Pair
Shortest Path Problem} (SPSPP) dans la littérature algorithmique. C'est
exactement le problème que résout un service de navigation type Google
Maps~: minimiser non pas la distance, mais le \textbf{temps total de
trajet}.

\subsection{Distance Euclidienne Sphérique~: la Formule de Haversine}
\label{ssec:haversine}

Avant même de raisonner sur le graphe routier, la plateforme a besoin de
mesurer une distance \og à vol d'oiseau \fg{} entre deux points GPS --- pour
un affichage rapide de proximité, et comme composante de l'heuristique de
l'algorithme A\* (section~\ref{ssec:dijkstra-astar}). La Terre étant
assimilable à une sphère, une simple différence de coordonnées cartésiennes
serait erronée~; on utilise la \textbf{formule de Haversine}.

Pour deux points de latitude et longitude $(\varphi_1, \lambda_1)$ et
$(\varphi_2, \lambda_2)$, exprimés en radians~:

\begin{align}
  a &= \sin^2\!\left(\frac{\Delta\varphi}{2}\right)
       + \cos(\varphi_1)\cos(\varphi_2)\sin^2\!\left(\frac{\Delta\lambda}{2}\right) \label{eq:haversine-a}\\
  c &= 2 \cdot \operatorname{atan2}\!\left(\sqrt{a},\, \sqrt{1-a}\right) \label{eq:haversine-c}\\
  d &= R \cdot c \label{eq:haversine-d}
\end{align}

où $\Delta\varphi = \varphi_2 - \varphi_1$, $\Delta\lambda = \lambda_2 -
\lambda_1$, et $R = 6\,371$~km est le rayon moyen terrestre. La distance $d$
obtenue est une \textbf{borne inférieure} de la distance réelle sur route~:
c'est cette propriété qui la rend utilisable comme heuristique admissible
pour A\* (section~\ref{ssec:astar-heuristique}).

\subsection{Algorithmes de Recherche du Plus Court Chemin}
\label{ssec:dijkstra-astar}

\subsubsection{Algorithme de Dijkstra}

L'algorithme de \textbf{Dijkstra}~[16], publié en 1959, résout le
problème~\eqref{eq:plus-court-chemin} pour un graphe à poids positifs. Il
maintient, pour chaque sommet $v$, une
distance provisoire $\delta(v)$ --- le meilleur coût connu pour atteindre $v$
depuis $s$ --- initialisée à $\delta(s) = 0$ et $\delta(v) = +\infty$ pour
tout $v \neq s$. À chaque itération, l'algorithme extrait d'une file de
priorité le sommet $u$ non visité de plus petite distance $\delta(u)$, puis
\textbf{relâche} chacune de ses arêtes sortantes~: pour tout voisin $v$ tel
que $(u, v) \in E$,

\begin{equation}
  \text{si } \delta(u) + w(u, v) < \delta(v) \text{, alors }
  \delta(v) \leftarrow \delta(u) + w(u, v) \text{ et } \pi(v) \leftarrow u
  \label{eq:relaxation}
\end{equation}

où $\pi(v)$ mémorise le prédécesseur de $v$ sur le meilleur chemin trouvé.
L'algorithme se termine lorsque $t$ est extrait de la file~; le chemin optimal
se reconstruit en remontant $\pi(t), \pi(\pi(t)), \dots$ jusqu'à $s$.

\textbf{Correction de l'algorithme.} La preuve de correction repose sur
l'invariant suivant~: lorsqu'un sommet $u$ est extrait de la file avec la
distance minimale parmi les sommets non visités, $\delta(u)$ est
\textbf{définitivement} optimale. En effet, tout chemin alternatif vers $u$
devrait nécessairement passer par un sommet encore non visité $u'$ avec
$\delta(u') \geq \delta(u)$~; comme tous les poids sont strictement positifs
(hypothèse de la section~\ref{ssec:modele-graphe}), ce chemin ne pourrait pas
être plus court. Cette propriété échoue en présence de poids négatifs, ce qui
justifie l'exigence $w(u,v) > 0$ posée en~\eqref{eq:poids-arete}.

\textbf{Complexité.} Avec une file de priorité implémentée par un tas
binaire, la complexité de Dijkstra est~:

\begin{equation}
  O\big((|V| + |E|) \log |V|\big)
  \label{eq:complexite-dijkstra}
\end{equation}

Sur un graphe représentant le réseau routier d'un pays entier, $|V|$ se
compte en millions de sommets~: appliqué naïvement à chaque requête,
Dijkstra est trop lent pour un usage interactif, d'où le recours à
l'heuristique A\* puis aux pré-calculs décrits ci-après.

\subsubsection{Algorithme A\* et heuristique admissible}
\label{ssec:astar-heuristique}

L'algorithme \textbf{A\*}~[17] (A-star), publié en 1968, améliore Dijkstra en orientant la
recherche vers la destination $t$ au moyen d'une \textbf{heuristique}
$h : V \rightarrow \mathbb{R}^{+}$. Pour chaque sommet $v$ exploré, A\*
calcule~:

\begin{equation}
  f(v) = g(v) + h(v)
  \label{eq:astar-f}
\end{equation}

où $g(v)$ est le coût réel déjà parcouru depuis $s$ jusqu'à $v$ (identique au
$\delta(v)$ de Dijkstra), et $h(v)$ est une \textbf{estimation optimiste} du
coût restant de $v$ à $t$. Contrairement à Dijkstra, qui explore les sommets
par ordre croissant de $g(v)$ dans toutes les directions, A\* explore par
ordre croissant de $f(v)$~: il privilégie les sommets qui semblent
\og prometteurs \fg{} vis-à-vis de la destination.

\textbf{Choix de l'heuristique.} Dans notre contexte, $h(v)$ est la distance
de Haversine (équation~\eqref{eq:haversine-d}) entre $v$ et $t$, divisée par
la vitesse maximale autorisée pour le mode choisi~:

\begin{equation}
  h(v) = \frac{d_{\text{Haversine}}(v, t)}{v_{\text{max}}(m)}
  \label{eq:heuristique-h}
\end{equation}

Cette heuristique est \textbf{admissible}~: elle ne surestime jamais le coût
réel restant, puisque la distance en ligne droite est, par construction,
inférieure ou égale à la distance réelle sur route
($d_{\text{Haversine}}(v,t) \leq d_{\text{route}}(v,t)$). Cette propriété
d'admissibilité \textbf{garantit} que A\* retourne toujours le chemin
optimal, exactement comme Dijkstra, mais en explorant un nombre de sommets
significativement inférieur~: intuitivement, Dijkstra explore un disque
centré sur $s$, tandis que A\* explore une ellipse étirée en direction de
$t$.

\subsubsection{Accélération par hiérarchies de contraction}

Même A\* reste trop lent pour répondre en quelques millisecondes sur un
graphe de plusieurs millions de sommets. Les moteurs de routing en production
(dont Valhalla, utilisé par la plateforme --- section~\ref{ssec:justification-algo})
ajoutent un \textbf{pré-calcul hors-ligne} appelé \textbf{Contraction
Hierarchies} (hiérarchies de contraction)~: le graphe est prétraité une fois
pour toutes en classant les arêtes par importance (autoroute $>$ route
nationale $>$ voie locale) et en insérant des arêtes de raccourci
(\textit{shortcuts}) qui résument de longs segments de route principale en
une seule arête équivalente. À la requête, le moteur exécute un \textbf{A\*
bidirectionnel}~: deux recherches simultanées, l'une progressant depuis $s$,
l'autre depuis $t$, qui se rencontrent au milieu du graphe enrichi de
raccourcis. Cette combinaison ramène le temps de réponse de plusieurs
secondes à quelques millisecondes, pour un résultat toujours optimal.

\medskip
\noindent\textbf{Synthèse.} Le tableau ci-dessous récapitule la correspondance
entre les concepts mathématiques formalisés dans cette section et leur rôle
concret dans la plateforme.

\begin{table}[H]
\centering
\caption{Correspondance entre concepts mathématiques et éléments applicatifs}
\label{tab:synthese-math}
\begin{tabular}{>{\bfseries}m{4.2cm} m{9.5cm}}
\toprule
\rowcolor{eecgreen}
\multicolumn{1}{>{\bfseries\color{white}}m{4.2cm}}{Concept} &
\multicolumn{1}{>{\bfseries\color{white}}m{9.5cm}}{Rôle dans la plateforme} \\
\midrule
\rowcolor{eeclightgreen!60}
Graphe $G=(V,E,w)$ & Réseau routier camerounais (données OpenStreetMap) \\
Sommet $v \in V$ & Intersection / carrefour routier \\
\rowcolor{eeclightgreen!60}
Arête $(u,v) \in E$ & Tronçon de route entre deux intersections \\
Poids $w(u,v)$ & Temps de parcours du tronçon, selon le mode choisi \\
\rowcolor{eeclightgreen!60}
Plus court chemin & Itinéraire optimal entre point de départ et entité EEC \\
Dijkstra & Algorithme de base garantissant l'optimalité \\
\rowcolor{eeclightgreen!60}
A\* et heuristique $h$ & Recherche guidée, accélérée par la distance de Haversine \\
Contraction Hierarchies & Pré-calcul assurant une réponse en temps réel \\
\bottomrule
\end{tabular}
\end{table}

%-------------

\medskip
\noindent\textbf{Bilan du chapitre.}
Ce chapitre a formalisé l'ensemble des besoins de la plateforme en
vingt-quatre exigences fonctionnelles et douze exigences non-fonctionnelles,
couvrant les domaines de l'authentification, de la cartographie, de
l'administration des entités, de l'export, des statistiques, de la
traçabilité, de la gestion hiérarchique des comptes et des services du
visiteur authentifié --- dont l'itinéraire réel et la navigation~3D. Les six modèles UML
d'analyse --- contexte (six acteurs), packages, cas d'utilisation
(UC01--UC29), classes métier et activité --- ont défini avec précision les
acteurs, les modules, les interactions et les entités du domaine institutionnel.
La section de conception a complété cette analyse en présentant l'architecture
technique à cinq services Docker, le modèle de données PostGIS (quinze tables),
le diagramme de séquence d'authentification sécurisée et le diagramme de
classes technique des modèles Django. Enfin, la section de modélisation
mathématique a formalisé le graphe pondéré du réseau routier, la formule de
Haversine et les algorithmes Dijkstra et A\* qui fondent scientifiquement le
service de calcul d'itinéraire de la plateforme.
Le \textbf{Chapitre~3} présentera l'environnement de développement retenu,
justifiera les choix technologiques et exposera les résultats concrets de
l'implémentation.

\clearpage

% ============================================================
%  CHAPITRE 3 — IMPLÉMENTATION ET RÉSULTATS
% ============================================================
\chapter{Implémentation et Résultats}
\label{chap:implementation}

\begin{chapplan}
  \item \textbf{Section~\ref{sec:outils} --- Outils et Technologies}~:
    justification des choix technologiques pour chaque composant de la
    pile logicielle~; algorithmes et moteur de calcul d'itinéraire
    (OSRM, Valhalla, décodage de polyligne)~; tableau récapitulatif des
    outils, versions et rôles.
  \item \textbf{Section~\ref{sec:deploiement} --- Diagramme de Déploiement}~:
    distribution des cinq services sur les nœuds physiques et virtuels~;
    configuration réseau Docker Compose.
  \item \textbf{Section~\ref{sec:resultats} --- Résultats}~:
    présentation annotée des interfaces fonctionnelles~: carte interactive,
    authentification, tableaux de bord par rôle, export, journal d'audit,
    itinéraire réel et navigation~3D~; confrontation aux exigences du
    Chapitre~2.
  \item \textbf{Section~\ref{sec:discussion} --- Discussion}~:
    vérification des objectifs, analyse critique des forces et des limites
    de la solution réalisée.
\end{chapplan}

\lettrine[lines=2, loversize=0.05]{L}{es} décisions de conception
formalisées au Chapitre~2 se concrétisent dans ce chapitre en une
implémentation fonctionnelle et déployable. La
section~\ref{sec:outils} justifie les choix technologiques opérés pour
chaque composant de la pile logicielle et dresse un tableau récapitulatif
des outils utilisés. La section~\ref{sec:deploiement} présente le diagramme
de déploiement UML qui décrit la distribution physique des services.
La section~\ref{sec:resultats} expose les résultats à travers les
interfaces fonctionnelles de la plateforme, commentées au regard des
exigences définies au Chapitre~2. La section~\ref{sec:discussion} propose
enfin une analyse critique des forces et des limites de la solution.

%------------------------------------------------------------
\section{Outils et Technologies}
\label{sec:outils}
%------------------------------------------------------------

\subsection{Justification des Choix Technologiques}
\label{ssec:justification}

Le choix de chaque technologie résulte d'une évaluation explicite des
alternatives disponibles au regard des contraintes identifiées~: contraintes
institutionnelles (déploiement sur infrastructure limitée, données
sensibles), contraintes fonctionnelles (requêtes spatiales, RBAC
hiérarchique, publication OGC) et contraintes de maintenabilité
(logiciels libres, documentation active, communauté de support).

\subsubsection{Django~5 et Django REST Framework}

Django~5~[18] a été retenu comme framework backend principal
pour trois raisons déterminantes. En premier lieu, son extension spatiale
native GeoDjango offre une intégration directe avec PostGIS~: les requêtes
de distance (\texttt{Distance()}), les filtres spatiaux (\texttt{contains},
\texttt{intersects}) et les champs de géométrie (\texttt{PointField},
\texttt{MultiPolygonField}) sont disponibles via l'ORM sans SQL manuel.
En second lieu, le module d'authentification de Django~5 propose une gestion
de sessions robuste par cookies \textit{HttpOnly}~; ce mécanisme est
nativement protégé contre les attaques de type XSS, contrairement aux
jetons JWT transmis dans le stockage local du navigateur, qui présentaient
une vulnérabilité documentée dans la version antérieure de la plateforme.
En troisième lieu, l'écosystème Python disponible ---
\texttt{openpyxl} pour la lecture des fichiers Excel EEC,
\texttt{reportlab} pour la génération PDF, \texttt{django-cors-headers}
pour la politique d'origines croisées --- couvre l'ensemble des besoins
fonctionnels sans dépendances tierces lourdes. Django REST
Framework~[14] complète Django en exposant les ressources
institutionnelles sous forme d'API REST JSON consommée par le frontend.

\subsubsection{PostgreSQL~16 et PostGIS~3.4}

PostgreSQL~16 associé à PostGIS~3.4~[12] constitue le
standard de facto pour les bases de données spatiales open source~: il est
utilisé dans des projets cartographiques institutionnels à l'échelle nationale
et internationale, dispose d'une indexation spatiale performante
(\textit{R-tree GiST}) et garantit la conformité aux spécifications OGC
Simple Features. La fonction \texttt{ST\_Distance()} native permet
d'identifier la paroisse la plus proche d'une position GPS en une seule
requête SQL, sans traitement applicatif. Aucun SGBD alternatif ne combine
ces fonctionnalités spatiales et une intégration aussi mature avec GeoDjango.

\subsubsection{GeoServer~2.26}

GeoServer~2.26~[19] a été retenu pour la publication des
couches cartographiques institutionnelles au regard de sa conformité native
aux protocoles OGC WMS~1.3.0 et WFS~2.0.0, qui garantissent
l'interopérabilité avec tout client cartographique standard. La
configuration des styles visuels via des fichiers SLD permet de définir les
couleurs institutionnelles de l'EEC (vert synodal, fond clair) directement
sur le serveur, sans traitement côté client. L'alternative d'une
publication de fichiers GeoJSON statiques, utilisée dans la version
antérieure, était incompatible avec la mise à jour dynamique des couches
et ne répondait pas aux exigences ENF05 d'interopérabilité OGC.

\subsubsection{Next.js~16}

Next.js~16~[20] avec l'App Router a été retenu pour le frontend
pour sa capacité à combiner rendu côté serveur (\textit{SSR}) et rendu côté
client (\textit{CSR}) dans la même application. Le \textit{SSR} est
utilisé pour les pages d'administration~: les données sensibles ne sont
jamais envoyées au navigateur avant vérification de la session, ce qui
réduit la surface d'attaque. Le rendu côté client est utilisé pour la carte
interactive, qui requiert des interactions dynamiques non compatibles avec
le \textit{SSR}. L'App Router de Next.js~16 permet de définir des
\textit{layouts} partagés par rôle d'administration (SUPER/REGION,
DISTRICT, PAROISSE), ce qui simplifie l'implémentation du RBAC côté
frontend.

\subsubsection{Leaflet.js et MarkerCluster}

Leaflet.js~[21] a été retenu pour la bibliothèque
cartographique cliente pour sa légèreté (environ 42~Ko minifié), sa licence
libre et l'étendue de son écosystème de plugins. Le plugin
\texttt{Leaflet.MarkerCluster} implémente le regroupement automatique des
858~marqueurs de paroisses en indicateurs numériques~: sans ce mécanisme,
la densité des marqueurs rendrait la carte illisible aux niveaux de zoom
nationaux. La prise en charge native des couches \textit{TileLayer.WMS}
permet de consommer directement les couches GeoServer sans adaptation.

\subsubsection{MapLibre~GL~JS}
\label{ssec:justification-maplibre}

\textbf{MapLibre~GL~JS}~[25] a été retenu pour assurer la vue de navigation
guidée en trois dimensions (section~\ref{ssec:res-avance}), en complément de
Leaflet qui reste la bibliothèque cartographique principale de
l'application. Ce choix répond à un besoin que Leaflet ne couvre pas
nativement~: le rendu vectoriel accéléré par \textbf{WebGL}, seul capable
d'afficher une carte inclinée ($\textit{pitch}$) et pivotée dynamiquement
selon le cap de déplacement de l'utilisateur, à la manière d'une application
de navigation grand public. Trois raisons ont motivé ce choix plutôt qu'une
alternative propriétaire (Google Maps, Mapbox GL)~: une licence entièrement
libre et sans clé d'API~; la compatibilité directe avec les styles
vectoriels ouverts (\textit{OpenFreeMap}), évitant toute dépendance à un
fournisseur commercial~; et une empreinte limitée sur l'application
principale, puisque le composant n'est chargé qu'à la demande
(\textit{lazy loading}), au déclenchement d'une navigation, puis démonté
immédiatement à son arrêt. La carte publique et les interfaces
d'administration continuent ainsi de reposer exclusivement sur Leaflet~1.9,
plus léger (42~Ko minifié) et suffisant pour un usage non immersif.

\subsubsection{Redis~7 et Docker Compose}

Redis~7~[22] assure deux fonctions complémentaires~:
le stockage des sessions Django côté serveur (évitant leur persistance
en base de données relationnelle) et la mise en cache des réponses API
fréquentes. Docker Compose~[15] orchestre l'ensemble des
cinq services dans une configuration déclarative reproductible sur toute
machine hôte disposant de Docker~: c'est la réponse directe à l'exigence
ENF06 de portabilité et à la limite de déploiement identifiée dans la
plateforme GeoEEC~v1.0.

%-----
\subsection{Algorithmes et Moteur de Calcul d'Itinéraire}
\label{ssec:justification-algo}

La section~\ref{sec:modelisation-math} a formalisé mathématiquement le
problème du plus court chemin sur un graphe routier. Cette sous-section
détaille l'implémentation concrète de cette modélisation~: le moteur
utilisé, la stratégie de résilience à trois niveaux, et l'algorithme de
décodage des tracés.

\subsubsection{Choix du moteur de routing~: Valhalla et OSRM}

Construire et maintenir un graphe routier national --- avec ses millions de
sommets, ses sens interdits, ses limitations de vitesse et ses mises à jour
continues --- représente un travail d'infrastructure considérable, hors du
périmètre raisonnable d'un projet institutionnel. La plateforme délègue donc
le calcul brut du plus court chemin à des \textbf{moteurs de routing
externes, gratuits et open source}, construits sur les données
\textbf{OpenStreetMap}~:
\begin{itemize}[itemsep=2pt, topsep=4pt, leftmargin=1.5cm]
  \item \textbf{OSRM}~[23] (\textit{Open Source Routing Machine}) --- instance
    publique de démonstration --- est sollicité en premier~: il implémente
    l'algorithme de Contraction Hierarchies décrit en
    section~\ref{ssec:dijkstra-astar} et répond en quelques centaines de
    millisecondes pour le mode voiture~;
  \item \textbf{Valhalla}~[24] (instance publique FOSSGIS) est sollicité en
    second recours~: il est \textbf{multimodal} (voiture, vélo, piéton) et
    implémente la même famille d'algorithmes --- A\* bidirectionnel
    accéléré par hiérarchies --- avec un support natif des trois profils de
    vitesse nécessaires à la plateforme~;
  \item en \textbf{dernier recours}, si aucun des deux services n'est
    joignable, la plateforme retombe sur un calcul de distance directe
    (formule de Haversine, équations~\eqref{eq:haversine-a}
    à~\eqref{eq:haversine-d}), garantissant que le service reste toujours
    fonctionnel, avec un tracé rectiligne explicitement signalé comme tel à
    l'utilisateur.
\end{itemize}

Cette architecture à trois niveaux de repli (\textit{fallback en cascade})
est implémentée dans le module isolé \texttt{backend/apps/visitors/routing.py}~:
si l'institution héberge un jour sa propre instance Valhalla ou OSRM, seule
l'URL du service change --- aucune autre partie de l'application n'est
affectée.

\subsubsection{Algorithme de sélection du fournisseur}

La fonction \texttt{calculer\_route()} implémente l'algorithme de
résilience suivant~:

\begin{enumerate}[itemsep=2pt, topsep=4pt, leftmargin=1.5cm]
  \item tenter le calcul via OSRM (profil voiture)~;
  \item en cas d'échec (service injoignable, délai dépassé, aucun itinéraire
    trouvé), tenter Valhalla avec le \textit{costing} correspondant au mode
    demandé (\texttt{auto}, \texttt{bicycle} ou \texttt{pedestrian})~;
  \item en cas de double échec, calculer un trajet direct par la formule de
    Haversine, avec une vitesse moyenne conventionnelle par mode
    (40~km/h en voiture par défaut, 15~km/h à vélo, 5~km/h à pied)~;
  \item dans tous les cas, retourner une structure de données uniforme~:
    mode, distance en kilomètres, durée en minutes, géométrie du tracé
    (liste de points \texttt{[latitude, longitude]}) et liste d'instructions
    virage-par-virage.
\end{enumerate}

Cette fonction ne lève \textbf{jamais} d'exception vers l'appelant~: la
disponibilité du service de calcul d'itinéraire prime sur l'exactitude
optimale du tracé, conformément à l'exigence ENF03 de disponibilité
(section~\ref{ssec:exigences}).

\subsubsection{Estimation de la durée par mode}

Lorsque le tracé provient d'OSRM (profil voiture uniquement), la durée pour
les modes vélo et piéton n'est pas fournie par le service~: elle est
recalculée à partir de la distance du tracé routier suivi et d'une vitesse
moyenne conventionnelle par mode~:

\begin{equation}
  \text{durée}_{\min} = \frac{\text{distance}_{\text{km}}}{v_{\text{mode}}} \times 60
  \label{eq:duree-mode}
\end{equation}

avec $v_{\text{bicycle}} = 15$~km/h et $v_{\text{pedestrian}} = 5$~km/h. Le
tracé géographique suivi reste néanmoins celui du réseau routier réel~;
seule l'estimation temporelle est ajustée au mode, une approximation
acceptée en l'absence d'un service multimodal unique répondant de façon
fiable pour les trois profils.

\subsubsection{Décodage de la polyligne encodée}

Les moteurs de routing ne transmettent pas les centaines de points GPS d'un
tracé en clair~: ils les compressent au format \textit{encoded polyline}.
L'algorithme de décodage, implémenté dans \texttt{decode\_polyline()},
repose sur trois principes~:
\begin{enumerate}[itemsep=2pt, topsep=4pt, leftmargin=1.5cm]
  \item chaque coordonnée n'est pas stockée en valeur absolue mais en
    \textbf{delta} par rapport au point précédent --- les écarts entre
    points consécutifs d'un tracé routier étant faibles, ils se compressent
    fortement~;
  \item chaque delta est encodé en base~64 modifiée, par blocs successifs de
    5~bits, avec un bit de continuation indiquant si le nombre se poursuit
    sur le caractère suivant~;
  \item la précision est fixée à $10^{-6}$ degré~: chaque entier décodé est
    divisé par $1\,000\,000$ pour retrouver la coordonnée décimale.
\end{enumerate}
Le décodage reconstruit ainsi la liste de points
\texttt{[[latitude, longitude], \ldots]} directement exploitable par la
bibliothèque cartographique Leaflet côté frontend, sans transformation
supplémentaire.

%-----
\subsection{Récapitulatif des Outils Utilisés}
\label{ssec:outils-table}

Le Tableau~\ref{tab:outils} dresse la liste complète des outils et
technologies utilisés dans le projet, avec leur version exacte et leur
rôle précis.

\begin{table}[H]
\centering
\caption{Récapitulatif des outils et technologies de la plateforme EEC}
\label{tab:outils}
\setlength{\tabcolsep}{5pt}
\renewcommand{\arraystretch}{1.22}
\resizebox{\linewidth}{!}{%
\begin{tabular}{%
  >{\bfseries}m{3.6cm}
  >{\centering\arraybackslash}m{1.6cm}
  m{9.0cm}}
\toprule
\rowcolor{eecgreen}
\multicolumn{1}{>{\bfseries\color{white}}m{3.6cm}}{Outil / Technologie} &
\multicolumn{1}{>{\bfseries\color{white}\centering\arraybackslash}m{1.6cm}}{Version} &
\multicolumn{1}{>{\bfseries\color{white}}m{9.0cm}}{Rôle dans le projet} \\
\midrule
\rowcolor{eeclightgreen!60}
Python           & 3.12   & Langage principal du backend \\
Django           & 5.0    & Framework web backend~; ORM~; middleware RBAC~; sessions \\
\rowcolor{eeclightgreen!60}
GeoDjango        & 5.0    & Extension spatiale de Django~; interface avec PostGIS \\
Django REST Framework & 3.15 & Construction et sérialisation des endpoints API REST \\
\rowcolor{eeclightgreen!60}
PostgreSQL       & 16     & Système de gestion de base de données relationnelle \\
PostGIS          & 3.4    & Extension spatiale~; types géométriques~; requêtes de distance \\
\rowcolor{eeclightgreen!60}
GeoServer        & 2.26   & Serveur cartographique OGC~; publication WMS et WFS~; styles SLD \\
Redis            & 7      & Cache serveur~; stockage des sessions Django \\
\rowcolor{eeclightgreen!60}
Next.js          & 16     & Framework frontend React~; App Router~; SSR et CSR \\
TypeScript       & 5.x    & Typage statique du code frontend \\
\rowcolor{eeclightgreen!60}
Tailwind CSS     & 3.x    & Framework de mise en forme CSS utilitaire \\
Leaflet.js       & 1.9    & Bibliothèque cartographique JavaScript interactive \\
\rowcolor{eeclightgreen!60}
Leaflet.MarkerCluster & 1.5 & Regroupement automatique des marqueurs sur la carte \\
openpyxl         & 3.1    & Lecture et écriture des fichiers Excel (.xlsx) \\
\rowcolor{eeclightgreen!60}
reportlab        & 4.x    & Génération des documents PDF (exports, rapports) \\
Docker           & 24     & Conteneurisation des services de la plateforme \\
\rowcolor{eeclightgreen!60}
Docker Compose   & 2.x    & Orchestration des cinq services en environnement multi-conteneurs \\
MapLibre GL JS & 5.24 & Rendu cartographique 3D (WebGL) pour la navigation guidée \\
\rowcolor{eeclightgreen!60}
OSRM / Valhalla & --- & Services de routing externes~; calcul d'itinéraire réel sur route \\
draw.io / diagrams.net & --- & Conception des diagrammes UML du rapport \\
\bottomrule
\end{tabular}}
\end{table}
%------------------------------------------------------------
\subsection{Analyse Comparative des Alternatives Technologiques}
\label{ssec:comparaison-tech}
%------------------------------------------------------------

Avant de retenir la pile logicielle finale, quatre domaines technologiques
clés ont été analysés~: le framework backend, la base de données spatiale,
le serveur cartographique et la bibliothèque cartographique cliente. Pour
chaque domaine, les alternatives les plus répandues ont été évaluées selon
des critères directement issus des exigences fonctionnelles et
non-fonctionnelles du Chapitre~2. Les symboles employés sont les suivants~:
\oui~critère entièrement satisfait, \non~critère non satisfait,
\partiel~critère partiellement satisfait.

% ---- Table 1 : Framework backend ----
\vspace{6pt}
\noindent Le Tableau~\ref{tab:comp-backend} compare les quatre frameworks
backend candidats selon les exigences de l'EEC~: support spatial natif,
gestion sécurisée des sessions, middleware RBAC configurable,
écosystème compatible avec les formats de données institutionnels
(Excel, PDF) et licence libre.

\begin{table}[H]
\centering
\caption{Comparaison des frameworks backend candidats}
\label{tab:comp-backend}
\setlength{\tabcolsep}{5pt}
\renewcommand{\arraystretch}{1.22}
\resizebox{\linewidth}{!}{%
\begin{tabular}{%
  >{\bfseries}m{5.0cm}
  >{\centering\arraybackslash}m{2.2cm}
  >{\centering\arraybackslash}m{2.0cm}
  >{\centering\arraybackslash}m{2.0cm}
  >{\centering\arraybackslash}m{2.2cm}}
\toprule
\rowcolor{eecgreen}
\multicolumn{1}{>{\bfseries\color{white}}m{5.0cm}}{Critère} &
\multicolumn{1}{>{\bfseries\color{white}\centering\arraybackslash}m{2.2cm}}{\textbf{Django~5} \newline \small(retenu)} &
\multicolumn{1}{>{\bfseries\color{white}\centering\arraybackslash}m{2.0cm}}{Flask \newline \small 3.x} &
\multicolumn{1}{>{\bfseries\color{white}\centering\arraybackslash}m{2.0cm}}{FastAPI \newline \small 0.x} &
\multicolumn{1}{>{\bfseries\color{white}\centering\arraybackslash}m{2.2cm}}{Express.js \newline \small(Node)} \\
\midrule
Extension spatiale native (GeoDjango)
  & \oui & \non & \non & \non \\
\rowcolor{eeclightgreen!60}
Sessions \textit{HttpOnly} intégrées
  & \oui & \partiel & \non & \non \\
Middleware RBAC configurable
  & \oui & \partiel & \partiel & \partiel \\
\rowcolor{eeclightgreen!60}
Écosystème Python (openpyxl, reportlab)
  & \oui & \oui & \oui & \non \\
ORM intégré (requêtes spatiales)
  & \oui & \non & \non & \non \\
\rowcolor{eeclightgreen!60}
Protection CSRF native
  & \oui & \partiel & \non & \non \\
Journal d'audit via signaux intégrés
  & \oui & \non & \non & \non \\
\rowcolor{eeclightgreen!60}
Licence libre
  & \oui & \oui & \oui & \oui \\
\midrule
\rowcolor{eecgreen!20}
\textbf{Score (sur 8)}
  & \textbf{8/8} & 3/8 & 2/8 & 1/8 \\
\bottomrule
\end{tabular}}
\end{table}

% ---- Table 2 : Base de données spatiale ----
\vspace{6pt}
\noindent Le Tableau~\ref{tab:comp-bd} évalue les systèmes de gestion de
bases de données spatiales selon les besoins de géolocalisation et
d'intégration avec GeoDjango.

\begin{table}[H]
\centering
\caption{Comparaison des bases de données spatiales candidates}
\label{tab:comp-bd}
\setlength{\tabcolsep}{5pt}
\renewcommand{\arraystretch}{1.22}
\resizebox{\linewidth}{!}{%
\begin{tabular}{%
  >{\bfseries}m{5.5cm}
  >{\centering\arraybackslash}m{2.5cm}
  >{\centering\arraybackslash}m{2.0cm}
  >{\centering\arraybackslash}m{2.0cm}
  >{\centering\arraybackslash}m{2.2cm}}
\toprule
\rowcolor{eecgreen}
\multicolumn{1}{>{\bfseries\color{white}}m{5.5cm}}{Critère} &
\multicolumn{1}{>{\bfseries\color{white}\centering\arraybackslash}m{2.5cm}}{\textbf{PostgreSQL} \newline \textbf{+ PostGIS} \newline \small(retenu)} &
\multicolumn{1}{>{\bfseries\color{white}\centering\arraybackslash}m{2.0cm}}{MySQL \newline \small+ Spatial} &
\multicolumn{1}{>{\bfseries\color{white}\centering\arraybackslash}m{2.0cm}}{MongoDB \newline \small+ Geo} &
\multicolumn{1}{>{\bfseries\color{white}\centering\arraybackslash}m{2.2cm}}{SQLite \newline \small+ SpatiaLite} \\
\midrule
Types géométriques OGC complets
  & \oui & \partiel & \partiel & \partiel \\
\rowcolor{eeclightgreen!60}
Calcul de distance spatial natif
  & \oui & \partiel & \partiel & \partiel \\
Intégration native GeoDjango
  & \oui & \partiel & \non & \partiel \\
\rowcolor{eeclightgreen!60}
Indexation spatiale performante (GiST)
  & \oui & \partiel & \oui & \non \\
Connexion directe à GeoServer
  & \oui & \oui & \non & \non \\
\rowcolor{eeclightgreen!60}
Requêtes hiérarchiques complexes
  & \oui & \oui & \non & \oui \\
Standard de facto SIG open source
  & \oui & \non & \non & \non \\
\rowcolor{eeclightgreen!60}
Licence libre
  & \oui & \oui & \oui & \oui \\
\midrule
\rowcolor{eecgreen!20}
\textbf{Score (sur 8)}
  & \textbf{8/8} & 5/8 & 3/8 & 4/8 \\
\bottomrule
\end{tabular}}
\end{table}

% ---- Table 3 : Serveur cartographique ----
\vspace{6pt}
\noindent Le Tableau~\ref{tab:comp-carto} confronte les solutions de
publication cartographique selon les critères de conformité OGC, de
configuration de styles et de connexion à la base de données spatiale.

\begin{table}[H]
\centering
\caption{Comparaison des solutions de publication cartographique}
\label{tab:comp-carto}
\setlength{\tabcolsep}{5pt}
\renewcommand{\arraystretch}{1.22}
\resizebox{\linewidth}{!}{%
\begin{tabular}{%
  >{\bfseries}m{5.2cm}
  >{\centering\arraybackslash}m{2.3cm}
  >{\centering\arraybackslash}m{2.1cm}
  >{\centering\arraybackslash}m{2.1cm}
  >{\centering\arraybackslash}m{2.0cm}}
\toprule
\rowcolor{eecgreen}
\multicolumn{1}{>{\bfseries\color{white}}m{5.2cm}}{Critère} &
\multicolumn{1}{>{\bfseries\color{white}\centering\arraybackslash}m{2.3cm}}{\textbf{GeoServer} \newline \textbf{2.26} \newline \small(retenu)} &
\multicolumn{1}{>{\bfseries\color{white}\centering\arraybackslash}m{2.1cm}}{MapServer \newline \small 8.x} &
\multicolumn{1}{>{\bfseries\color{white}\centering\arraybackslash}m{2.1cm}}{TileServer\newline\small-GL} &
\multicolumn{1}{>{\bfseries\color{white}\centering\arraybackslash}m{2.0cm}}{GeoJSON \newline \small statique} \\
\midrule
Conformité WMS 1.3.0
  & \oui & \oui & \non & \non \\
\rowcolor{eeclightgreen!60}
Conformité WFS 2.0.0
  & \oui & \oui & \non & \non \\
Styles SLD configurables
  & \oui & \oui & \non & \non \\
\rowcolor{eeclightgreen!60}
Interface d'administration graphique
  & \oui & \non & \partiel & \non \\
Connexion directe à PostGIS
  & \oui & \oui & \non & \non \\
\rowcolor{eeclightgreen!60}
Mise à jour dynamique des couches
  & \oui & \oui & \partiel & \non \\
Licence libre
  & \oui & \oui & \oui & \oui \\
\rowcolor{eeclightgreen!60}
Adapté aux infrastructures limitées
  & \oui & \partiel & \oui & \oui \\
\midrule
\rowcolor{eecgreen!20}
\textbf{Score (sur 8)}
  & \textbf{8/8} & 6/8 & 3/8 & 2/8 \\
\bottomrule
\end{tabular}}
\end{table}

% ---- Table 4 : Bibliothèque cartographique ----
\vspace{6pt}
\noindent Le Tableau~\ref{tab:comp-leaflet} compare les bibliothèques
cartographiques JavaScript selon les critères de gratuité, de support des
couches WMS, de disponibilité d'un plugin de regroupement de marqueurs
et de légèreté.

\begin{table}[H]
\centering
\caption{Comparaison des bibliothèques cartographiques JavaScript}
\label{tab:comp-leaflet}
\setlength{\tabcolsep}{5pt}
\renewcommand{\arraystretch}{1.22}
\resizebox{\linewidth}{!}{%
\begin{tabular}{%
  >{\bfseries}m{5.2cm}
  >{\centering\arraybackslash}m{2.2cm}
  >{\centering\arraybackslash}m{2.1cm}
  >{\centering\arraybackslash}m{2.1cm}
  >{\centering\arraybackslash}m{2.1cm}}
\toprule
\rowcolor{eecgreen}
\multicolumn{1}{>{\bfseries\color{white}}m{5.2cm}}{Critère} &
\multicolumn{1}{>{\bfseries\color{white}\centering\arraybackslash}m{2.2cm}}{\textbf{Leaflet.js} \newline \textbf{1.9} \newline \small(retenu)} &
\multicolumn{1}{>{\bfseries\color{white}\centering\arraybackslash}m{2.1cm}}{OpenLayers \newline \small 9.x} &
\multicolumn{1}{>{\bfseries\color{white}\centering\arraybackslash}m{2.1cm}}{Google Maps \newline \small API} &
\multicolumn{1}{>{\bfseries\color{white}\centering\arraybackslash}m{2.1cm}}{Mapbox GL \newline \small JS} \\
\midrule
Licence libre, sans quota d'usage
  & \oui & \oui & \non & \non \\
\rowcolor{eeclightgreen!60}
Support natif couches WMS
  & \oui & \oui & \non & \partiel \\
Plugin MarkerCluster disponible
  & \oui & \partiel & \oui & \oui \\
\rowcolor{eeclightgreen!60}
Légèreté ($<$~100~Ko minifié)
  & \oui & \non & \non & \non \\
Couverture cartographique Cameroun
  & \oui & \oui & \oui & \oui \\
\rowcolor{eeclightgreen!60}
Gratuité totale (sans clé d'API)
  & \oui & \oui & \non & \non \\
Documentation et communauté actives
  & \oui & \oui & \oui & \oui \\
\rowcolor{eeclightgreen!60}
Adapté aux connexions à faible débit
  & \oui & \partiel & \non & \non \\
\midrule
\rowcolor{eecgreen!20}
\textbf{Score (sur 8)}
  & \textbf{8/8} & 6/8 & 3/8 & 4/8 \\
\bottomrule
\end{tabular}}
\end{table}

%------------------------------------------------------------
\section{Diagramme de Déploiement}
\label{sec:deploiement}
%------------------------------------------------------------

Le diagramme de déploiement (Figure~\ref{fig:diag-deploiement}) représente
la distribution physique et virtuelle des composants logiciels de la
plateforme sur les nœuds d'exécution. Deux nœuds principaux sont identifiés~:
le \textbf{serveur institutionnel} et le \textbf{navigateur client}.

Sur le serveur institutionnel, Docker Compose crée un réseau interne
(\texttt{eec\_network}) dans lequel les cinq conteneurs communiquent~: le
conteneur \texttt{backend} (Django~5, port~8000) expose l'API REST au
frontend et interroge \texttt{db} (PostgreSQL/PostGIS, port~5432) et
\texttt{cache} (Redis, port~6379)~; le conteneur \texttt{geoserver}
(port~8080) accède à \texttt{db} pour lire les données géographiques et
répond aux requêtes WMS/WFS~; le conteneur \texttt{frontend} (Next.js~16,
port~3000) sert l'interface web. Vers l'extérieur, seul le port~3000
(interface) est exposé~; les services internes restent cloisonnés dans
le réseau Docker. Les volumes persistants (\texttt{postgres\_data},
\texttt{geoserver\_data}) garantissent la conservation des données entre
les redémarrages de conteneurs.

Sur le navigateur client, Leaflet.js charge les tuiles de fond de carte
depuis OpenStreetMap et interroge directement GeoServer (port~8080) pour
les couches institutionnelles WMS/WFS~; les requêtes API REST transitent
vers le backend via le frontend Next.js.

\begin{figure}[H]
  \center
  \includegraphics[width=1\textwidth]{deploiement.png}
  \caption{Diagramme de déploiement UML — Distribution des services Docker
Compose}
  \label{fig:diag-contexte}
\end{figure}

%------------------------------------------------------------
\section{Résultats}
\label{sec:resultats}
%------------------------------------------------------------

Cette section présente les interfaces fonctionnelles de la plateforme
implémentée, commentées au regard des exigences définies au
Chapitre~2. Les captures d'écran sont organisées selon les quatre
grands modules~: interface publique cartographique, authentification,
administration multi-niveaux et gestion des données.

%-----
\subsection{Page d'Accueil (\textit{Landing Page})}
\label{ssec:res-landing}

La page d'accueil est la première page présentée à tout visiteur accédant à
la plateforme. Elle met en avant trois actions principales~: se connecter
(administrateurs et visiteurs enregistrés), s'inscrire (nouveau visiteur) et
explorer la carte publique sans authentification.

\begin{figure}[H]
  \center
  \includegraphics[width=1\textwidth]{Capture d'écran 2026-07-03 035902.png}
  \caption{Page d'accueil de la plateforme}
  \label{fig:res-landing}
\end{figure}

\subsection{Interface Publique Cartographique}
\label{ssec:res-carte}

L'interface publique est accessible sans authentification depuis la racine
de la plateforme. Elle répond directement aux exigences EF03, EF04, EF05,
EF15 et EF16 relatives à la cartographie interactive et à l'accès public.

La Figure~\ref{fig:res-carte-publique} illustre la carte au niveau de zoom
national~: les 858~paroisses importées sont regroupées en indicateurs de
cluster numériques selon leur densité géographique. Le fond de carte
OpenStreetMap fournit le contexte géographique camerounais~; les couches
institutionnelles (polygones des régions synodales, marqueurs des paroisses
et des œuvres) sont servies par GeoServer en WMS. La barre de filtres
hiérarchiques en haut de la carte permet de sélectionner une région synodale,
un district et un type d'entité~; la carte se met à jour dynamiquement
sans rechargement de page.

\begin{figure}[H]
  \center
  \includegraphics[width=1\textwidth]{Capture d'écran 2026-07-03 003106.png}
  \caption{Carte interactive publique --- vue nationale avec clustering}
  \label{fig:res-carte-publique}
\end{figure}

La Figure~\ref{fig:res-popup} montre le panneau de détail qui s'affiche lors
d'un clic sur un marqueur individuel~: nom de la paroisse, district et région
d'appartenance, coordonnées GPS, contacts et statut d'activité. Cette fiche
contextuelle répond à l'exigence EF06 et est accessible à tout utilisateur
sans connexion.

\begin{figure}[H]
  \center
  \includegraphics[width=1\textwidth]{indi.png}
  \caption{Panneau de détail d'une entité, affiché depuis la carte}
  \label{fig:res-popup}
\end{figure}

%-----
\subsection{Interface d'Authentification}
\label{ssec:res-auth}

La Figure~\ref{fig:res-login} présente la page de connexion de l'espace
d'administration (\texttt{/admin/login}). Le formulaire soumet les
identifiants à l'API backend~; en cas de succès, un cookie de session
\textit{HttpOnly} est transmis et l'utilisateur est redirigé vers le
tableau de bord correspondant à son rôle. En cas d'identifiants incorrects,
un message d'erreur est affiché sans révéler si l'identifiant ou le mot de
passe est la source de l'erreur, conformément aux bonnes pratiques de
sécurité recommandées par l'OWASP~[11]. Pour un compte dont
l'indicateur de première connexion est actif, la redirection s'effectue
automatiquement vers le formulaire de changement de mot de passe.
\begin{figure}[H]
  \center
  \includegraphics[width=1\textwidth]{Capture d'écran 2026-05-28 060539.png}
  \caption{Page de connexion de l'espace d'administration}
  \label{fig:res-login}
\end{figure}

 

%-----
\subsection{Interface d'Administration Multi-Niveaux}
\label{ssec:res-admin}

L'interface d'administration adapte son contenu et ses fonctionnalités au
rôle de l'utilisateur connecté, conformément aux exigences EF02 et EF07 à
EF09. Quatre interfaces distinctes sont générées selon le rôle~: SUPER/REGION,
DISTRICT et PAROISSE.

La Figure~\ref{fig:res-dashboard} illustre le tableau de bord de
l'administrateur national (rôle SUPER)~: il affiche les indicateurs de
synthèse nationaux (nombre total de paroisses actives, de districts, de
régions, d'ouvriers et d'œuvres), les statistiques de couverture GPS par
région et l'évolution des indicateurs sur les dernières années. Les
indicateurs sont calculés dynamiquement par le backend en filtrant les
données selon le périmètre du rôle~: un administrateur régional voit les
mêmes indicateurs mais limités à sa région synodale.

\begin{figure}[H]
  \center
  \includegraphics[width=1\textwidth]{Capture d'écran 2026-05-28 060821.png}
  \caption{Tableau de bord de l'administrateur national (rôle SUPER)}
  \label{fig:res-dashboard}
\end{figure}

La Figure~\ref{fig:res-paroisses} présente l'interface de gestion des
paroisses~: liste paginée filtrée par région et district, formulaire de
création et de modification avec validation des coordonnées GPS en temps
réel, et action de suppression avec confirmation. Le système de contrôle
d'accès RBAC restreint les opérations à la liste des paroisses appartenant
au périmètre synodal de l'utilisateur~: un administrateur de district ne
peut visualiser ni modifier les paroisses d'un autre district.

\begin{figure}[H]
  \center
  \includegraphics[width=1\textwidth]{paroisse.png}
  \caption{Interface d'administration --- Gestion des paroisses avec contrôle RBAC}
  \label{fig:res-paroisses}
\end{figure}

%-----
\subsection{Module d'Export}
\label{ssec:res-import}

\noindent\textit{Note de version~:} une version antérieure de la plateforme
proposait également un module d'import de fichiers Excel~; il a été retiré
du périmètre fonctionnel (voir section~\ref{sec:audit-donnees}) au profit
d'une saisie directement contrôlée dans l'interface d'administration, jugée
plus fiable au regard des incohérences relevées dans les fichiers sources.
Seul l'export subsiste dans la version actuelle.

Le module d'export (Figure~\ref{fig:res-import}) répond à l'exigence EF10~:
il permet de télécharger les listes d'entités --- paroisses, œuvres,
ouvriers --- filtrées selon le périmètre de l'utilisateur, au format Excel
(\texttt{.xlsx}) ou PDF, directement depuis l'interface d'administration.

\begin{figure}[H]
  \center
  \includegraphics[width=1\textwidth]{import.png}
  \caption{Module d'export --- Téléchargement des données au format Excel et PDF}
  \label{fig:res-import}
\end{figure}

%-----
\subsection{Journal d'Audit}
\label{ssec:res-audit}

Le journal d'audit (Figure~\ref{fig:res-audit}) est accessible
exclusivement à l'administrateur national. Il liste toutes les opérations
d'écriture effectuées sur les données institutionnelles, dans l'ordre
chronologique inversé~: création, modification ou suppression d'une entité,
avec l'auteur de l'action, le type d'entité concernée, l'identifiant de
l'objet et l'horodatage UTC. Des filtres par type d'action, type d'entité
et période permettent de cibler une recherche. En cliquant sur une entrée,
l'administrateur accède au détail de la modification~: valeurs avant et
après la mise à jour. Ce module répond aux exigences EF13 et EF14 relatives
à la traçabilité, et constitue l'un des apports distinctifs par rapport
à la plateforme GeoEEC~v1.0 qui ne disposait d'aucun mécanisme d'audit.

\begin{figure}[H]
  \center
  \includegraphics[width=1\textwidth]{test11.png}
  \caption{Journal d'audit --- Historique des opérations d'écriture sur les données}
  \label{fig:res-audit}
\end{figure}

%-----
\subsection{Fonctionnalités Avancées~: Itinéraire Réel et Navigation~3D}
\label{ssec:res-avance}

Au-delà des modules d'administration, la plateforme intègre deux
fonctionnalités destinées au visiteur authentifié, directement issues de la
modélisation mathématique du chapitre précédent (section~\ref{sec:modelisation-math})
et de son implémentation algorithmique (section~\ref{ssec:justification-algo}).

\begin{figure}[H]
  \center
  \includegraphics[width=1\textwidth]{Capture d'écran 2026-07-03 002903.png}
  \caption{Jtinéraire Réel}
  \label{fig:res-audit}
\end{figure}

\subsubsection{Calcul d'itinéraire réel sur le réseau routier}

L'onglet \og Parcours \fg{} permet au visiteur de définir un point de départ
--- sa position GPS (via \texttt{navigator.geolocation}) ou une entité de la
carte --- et une destination --- une paroisse ou une œuvre --- puis de
choisir un mode de déplacement (voiture, vélo, à pied). Le point d'entrée


\texttt{POST~/api/visitor/itineraires/route/} invoque la chaîne de
résilience OSRM~$\rightarrow$~Valhalla~$\rightarrow$~Haversine décrite en
section~\ref{ssec:justification-algo} et renvoie la distance, la durée, le
tracé géographique complet et les instructions virage-par-virage. Le
frontend dessine ce tracé sur la carte Leaflet~: un contour sombre
(\textit{casing}), une ligne colorée animée de pointillés défilants, des
marqueurs de départ et d'arrivée distincts, et un cadrage automatique de la
vue sur l'ensemble du trajet.

\begin{figure}[H]
  \center
  \includegraphics[width=1\textwidth]{Capture d'écran 2026-07-03 002823.png}
  \caption{Calcul d’itinéraire réel sur le réseau routier}
  \label{fig:res-audit}
\end{figure}

\subsubsection{Navigation~3D en temps réel}

Une fois l'itinéraire calculé, le visiteur peut démarrer une
\textbf{navigation guidée en trois dimensions}, comparable en expérience à
une application de navigation grand public. Cette vue est assurée par un
composant dédié, chargé paresseusement (\textit{lazy loading}) uniquement au
déclenchement de la navigation --- le reste de l'application demeure sur
Leaflet en 2D, conformément à l'exigence de légèreté du cahier des charges.
Elle repose sur \textbf{MapLibre~GL~JS}~[25] (rendu WebGL), qui permet une
inclinaison de caméra ($\textit{pitch} = 62°$) et une rotation dynamique
selon le cap de déplacement, à la différence de Leaflet qui ne rend qu'en
projection plane.

Trois mécanismes techniques assurent le réalisme du suivi~:
\begin{itemize}[itemsep=2pt, topsep=4pt, leftmargin=1.5cm]
  \item \textbf{Suivi GPS continu}~: l'API navigateur
    \texttt{watchPosition()} fournit la position en temps réel (haute
    précision, rafraîchissement sous la seconde)~;
  \item \textbf{Recalage sur trajectoire} (\textit{map matching})~: la
    position brute du GPS, sujette à des imprécisions, est projetée
    mathématiquement sur le segment du tracé le plus proche --- un problème
    de projection orthogonale d'un point sur un ensemble de segments
    consécutifs --- afin d'afficher un déplacement fluide et réaliste plutôt
    que \og bruité \fg{}~;
  \item \textbf{Cap dynamique}~: lorsque la vitesse mesurée dépasse un seuil
    significatif, le cap de la caméra suit le \textit{heading} du capteur
    GPS~; sinon, il retombe sur le relèvement (\textit{bearing}) calculé
    géométriquement entre les points consécutifs du tracé, évitant les
    rotations erratiques à l'arrêt.
\end{itemize}

Le temps et la distance restants sont recalculés à chaque position reçue~:
la distance restante est la différence entre la longueur totale du tracé et
la distance déjà parcourue le long de celui-ci (calculée par cumul des
segments Haversine), et la durée restante est réestimée à partir de la
vitesse instantanée du visiteur lorsqu'elle est disponible, avec un repli
sur une extrapolation proportionnelle à la distance restante sinon. L'arrivée
est détectée lorsque la distance restante descend sous un seuil de
30~mètres, ce qui déclenche un message de confirmation puis la fermeture
automatique de la vue~3D.

\begin{figure}[H]
  \center
  \includegraphics[width=1\textwidth]{Capture d'écran 2026-07-03 003407.png}
  \caption{demarage de nanvigation 3D}
  \label{fig:res-audit}
\end{figure}




Le choix technologique de MapLibre~GL~JS et son rôle dans l'architecture de
la plateforme sont justifiés en section~\ref{ssec:justification-maplibre}~;
le tableau~\ref{tab:outils} en récapitule la version et le rôle aux côtés
des autres outils du projet.

%------------------------------------------------------------
\section{Discussion}
\label{sec:discussion}
%------------------------------------------------------------

\subsection{Adéquation aux Objectifs}
\label{ssec:disc-objectifs}

Les six objectifs formulés en Introduction Générale sont confrontés aux
résultats obtenus. L'objectif~1 (modélisation du schéma de données
hiérarchique) est atteint~: le schéma PostGIS implémenté couvre les
22~régions synodales, 159~districts, 858~paroisses importées, 311~œuvres
et 685~ouvriers, avec toutes les contraintes d'intégrité référentielle et
les champs de géolocalisation. L'objectif~2 (API REST sécurisée avec RBAC)
est atteint~: l'API expose 47~endpoints documentés, protégés par le
middleware RBAC à quatre niveaux et par le journal d'audit automatique.
L'objectif~3 (configuration GeoServer OGC) est atteint~: trois couches
institutionnelles sont publiées en WMS~1.3.0 et WFS~2.0.0 avec styles SLD.
L'objectif~4 (interface cartographique Leaflet) est atteint~: la carte
interactive intègre le clustering automatique, les filtres hiérarchiques
et les popups de détail. L'objectif~5 (module d'administration multi-niveaux)
est atteint~: les quatre interfaces de rôle, les modules d'import Excel et
d'export PDF/Excel sont fonctionnels. L'objectif~6 (conteneurisation Docker)
est atteint~: la plateforme est entièrement déployable avec une commande
unique sur tout serveur disposant de Docker.

%-----
\subsection{Comparaison Chiffrée avec les Solutions Existantes}
\label{ssec:disc-comparatif}

Les résultats obtenus permettent à présent de confronter chiffrément notre
solution aux quatre solutions analysées au Chapitre~1 (section~\ref{sec:etat-art}).
Le tableau~\ref{tab:comparatif} reprend les onze critères jugés déterminants pour
une plateforme institutionnelle cartographique de l'EEC. Les symboles employés
sont les suivants~: \oui~critère entièrement satisfait, \non~critère non
satisfait, \partiel~critère partiellement satisfait, N/A~critère sans objet
pour cette catégorie de solution.

\begin{table}[H]
\centering
\caption{Tableau comparatif des solutions existantes et de la solution proposée}
\label{tab:comparatif}
\resizebox{\linewidth}{!}{%
\begin{tabular}{>{\bfseries}m{4.8cm}
                >{\centering\arraybackslash}m{2.0cm}
                >{\centering\arraybackslash}m{2.1cm}
                >{\centering\arraybackslash}m{2.1cm}
                >{\centering\arraybackslash}m{2.0cm}
                >{\centering\arraybackslash}m{2.2cm}}
\toprule
\textbf{Critère}
  & \textbf{GeoEEC v1.0}
  & \textbf{ArcGIS Online}
  & \textbf{Google Maps}
  & \textbf{QGIS / OL}
  & \textbf{Notre solution} \\
\midrule
Plateforme web multi-utilisateurs
  & \oui & \oui & \partiel & \non & \oui \\
\addlinespace[3pt]
Hiérarchie institutionnelle configurable
  & \partiel & \non & \non & \non & \oui \\
\addlinespace[3pt]
Authentification sécurisée (sessions)
  & \non & \oui & \oui & \non & \oui \\
\addlinespace[3pt]
RBAC à 4 niveaux synodaux
  & \partiel & \non & \non & \non & \oui \\
\addlinespace[3pt]
Serveur cartographique OGC (WMS/WFS)
  & \non & \oui & \non & \oui & \oui \\
\addlinespace[3pt]
Base de données spatiale (PostGIS)
  & \partiel & \oui & \oui & \oui & \oui \\
\addlinespace[3pt]
Journal d'audit des modifications
  & \non & \partiel & \non & \non & \oui \\
\addlinespace[3pt]
Import données Excel institutionnel
  & \oui & \non & \non & \non & \oui \\
\addlinespace[3pt]
Déploiement conteneurisé (Docker)
  & \non & N/A & N/A & \non & \oui \\
\addlinespace[3pt]
Logiciel libre / sans licence commerciale
  & \oui & \non & \non & \oui & \oui \\
\addlinespace[3pt]
Adapté au contexte institutionnel EEC
  & \partiel & \non & \non & \non & \oui \\
\bottomrule
\end{tabular}}
\end{table}

Comme au Chapitre~1, l'analyse comparative confirme, résultats à l'appui, que
notre solution est la seule à réunir simultanément l'ensemble des
caractéristiques requises~: elle ajoute désormais à cet acquis un calcul
d'itinéraire réel sur le réseau routier (section~\ref{ssec:disc-forces}) et une
navigation 3D, deux capacités qu'aucune des quatre solutions comparées
n'offre dans une configuration institutionnelle équivalente, gratuite et
souveraine.

%-----
\subsection{Forces de la Solution}
\label{ssec:disc-forces}

\textbf{Sécurité améliorée par rapport à la version antérieure.}
Le remplacement de l'authentification JWT par des sessions Django avec
cookies \textit{HttpOnly} et \textit{SameSite=Strict} élimine le vecteur
d'attaque XSS qui permettait la capture des jetons transmis dans le stockage
local du navigateur. La protection CSRF systématique sur toutes les
requêtes modifiantes renforce ce dispositif, conformément aux recommandations
de l'OWASP~[11].

\textbf{Interopérabilité cartographique.}
La publication des couches géographiques via GeoServer en WMS~1.3.0 et
WFS~2.0.0 rend les données de l'EEC consommables par tout client
cartographique compatible OGC~: QGIS, OpenLayers, ou un futur portail
institutionnel camerounais. Aucune solution concurrente analysée au
Chapitre~1 ne proposait cette combinaison.

\textbf{RBAC aligné sur la hiérarchie réelle de l'EEC.}
Le modèle de contrôle d'accès à quatre niveaux --- national, régional,
district et paroissial --- correspond exactement à la structure synodale
de l'institution~: un administrateur de district peut gérer les paroisses
de son périmètre sans accès aux données des districts voisins. Ce niveau
de granularité était absent de la version antérieure.

\textbf{Traçabilité complète et automatique.}
Tout enregistrement de modification, de création ou de suppression
déclenche automatiquement l'écriture d'une entrée dans le journal d'audit
via les signaux Django, sans intervention du développeur dans chaque vue.
Cette automatisation garantit l'exhaustivité de l'historique, impossible
à assurer manuellement.

\textbf{Déploiement reproductible.}
La configuration Docker Compose en cinq services garantit que l'environnement
d'exécution est identique en développement et en production~: plus aucune
divergence de versions de bibliothèques ou de configuration système ne peut
causer des défaillances à l'installation.

%-----
\subsection{Difficultés Rencontrées}
\label{ssec:disc-difficultes}

Si l'implémentation technique --- backend, cartographie, RBAC, algorithmes
de routing --- a représenté un volume de travail conséquent, ce n'est pas la
difficulté qui a le plus mis à l'épreuve la rigueur méthodologique de ce
projet. La phase d'\textbf{audit et de préparation des données}
(section~\ref{sec:audit-donnees}) a constitué, de loin, le défi le plus
exigeant~: contrairement à un développement logiciel classique où les
structures de données sont définies dès la conception, ce projet a dû
composer avec des données réelles, produites sur le terrain par des
services différents, à des dates différentes, sans référentiel commun.

Trois difficultés se sont révélées particulièrement coûteuses en temps
d'analyse. Premièrement, l'absence d'identifiant commun entre le fichier
d'enquête de géolocalisation et le référentiel officiel de catégorisation a
imposé le développement d'un algorithme de rapprochement flou
(\textit{fuzzy matching}) plutôt qu'une simple jointure~: sur les
516~paroisses de l'enquête, seules 267~présentaient un nom rigoureusement
identique dans les deux sources, la comparaison textuelle approchée ayant
permis d'en récupérer 66~supplémentaires. Deuxièmement, l'inversion des
colonnes de coordonnées GPS dans le fichier source --- \texttt{Coord\_x}
contenant en réalité la latitude --- est le type d'erreur silencieuse la
plus dangereuse~: elle ne provoque aucune erreur technique visible, mais
place chaque paroisse concernée à une position géographique incorrecte si
elle n'est pas détectée et corrigée en amont. Troisièmement, l'hétérogénéité
structurelle des fichiers (tableau croisé pour les œuvres, enregistrements
composites pour les ouvriers) a nécessité des traitements de restructuration
spécifiques à chaque source, aucune approche générique n'étant applicable.

Cette phase, bien qu'invisible dans l'interface finale de la plateforme,
constitue une part significative de la valeur ajoutée du projet~: une
plateforme cartographique n'a de valeur que si les données qu'elle expose
sont fiables, et c'est précisément ce travail de fiabilisation --- documenté,
outillé et reproductible plutôt qu'artisanal --- qui conditionne la
pertinence de l'ensemble des fonctionnalités présentées dans ce chapitre.

%-----
\subsection{Limites et Contraintes}
\label{ssec:disc-limites}

\textbf{Couverture GPS incomplète.}
L'état actuel de la base révèle que 480~paroisses sur 858, soit
56~\% de l'effectif total, ne disposent d'aucune coordonnée GPS valide.
Ces paroisses sont enregistrées en base mais ne s'affichent pas sur la
carte. La plateforme ne peut donc
pas prétendre à une couverture cartographique complète du réseau
ecclésiastique~; son utilité pour les décisions de planification
territoriale nationale reste limitée tant que ce taux n'est pas réduit.

\textbf{Dépendance à des services de routing externes.}
Le calcul d'itinéraire réel (section~\ref{ssec:justification-algo}) repose
sur des instances publiques et gratuites d'OSRM et de Valhalla, non
hébergées par l'institution. La plateforme reste fonctionnelle en cas
d'indisponibilité de ces services --- repli automatique sur la distance de
Haversine (section~\ref{ssec:haversine}) --- mais la qualité et la
disponibilité du tracé routier réel dépendent alors d'une infrastructure
tierce hors du contrôle de l'EEC, ce qui pose une question de souveraineté
et de continuité de service à long terme.

\textbf{Absence d'application mobile.}
La plateforme est conçue pour des écrans desktop et tablette
(résolution minimale 1\,024~$\times$~768~px, ENF09). Elle n'est pas
adaptée aux smartphones, qui constituent pourtant le principal outil
numérique des administrateurs paroissiaux en zone rurale camerounaise.
En l'absence d'application mobile dédiée, la collecte terrain des
coordonnées GPS des 480~paroisses non géolocalisées reste manuelle et
sans outillage dédié.

\textbf{Authentification à facteur unique.}
L'accès à l'espace d'administration repose sur un couple identifiant/mot
de passe sans mécanisme de double authentification (TOTP ou code
temporaire). Pour des comptes de niveau national disposant d'un accès
complet aux données de l'EEC, cette limite expose la plateforme à des
risques de compromission par capture ou réutilisation de mot de passe.

\medskip
\noindent\textbf{Bilan du chapitre.}
Ce chapitre a présenté et justifié les choix technologiques de la
plateforme, détaillé l'algorithmique du calcul d'itinéraire (moteurs OSRM et
Valhalla, décodage de polyligne), puis exposé les résultats à travers huit
interfaces fonctionnelles --- dont l'itinéraire réel et la navigation~3D ---
commentées au regard des exigences du Chapitre~2, et confronté chiffrément la
solution aux quatre solutions existantes. La discussion critique a
mis en évidence cinq forces distinctives --- sécurité des sessions, RBAC
synodal, interopérabilité OGC, audit automatique et déploiement conteneurisé
--- et quatre limites structurelles auxquelles la \textbf{Conclusion
Générale} propose des perspectives d'amélioration concrètes.

\clearpage

% ============================================================
%  CONCLUSION GÉNÉRALE
% ============================================================
\chapter*{Conclusion Générale}
\addcontentsline{toc}{chapter}{Conclusion Générale}
\markboth{Conclusion Générale}{Conclusion Générale}

\noindent
L'Église Évangélique du Cameroun gérait jusqu'alors ses 858~paroisses,
311~œuvres sociales et 685~ouvriers ecclésiastiques répartis sur
22~régions synodales et 159~districts sans aucun outil numérique
centralisé~: pas de visualisation géographique, pas de contrôle d'accès
formalisé selon les périmètres synodaux, pas de journal d'audit des
modifications, pas de publication cartographique interopérable. Ce travail
a abordé ce problème par la conception et le développement d'une
plateforme web cartographique institutionnelle, en s'appuyant sur
l'hypothèse qu'une architecture à services découplés, entièrement
conteneurisée, combinant un serveur cartographique OGC et un contrôle
d'accès hiérarchique aligné sur la structure synodale, permettrait d'y
répondre de façon complète et déployable dans un contexte institutionnel
à ressources limitées.

Cette hypothèse est validée par les résultats obtenus~: la plateforme
implémentée déploie en une seule commande cinq services
--- Django~5~+~GeoDjango, PostgreSQL~+~PostGIS~3.4, GeoServer~2.26,
Next.js~16 et Redis~7~--- et offre une carte interactive avec clustering
de 858~paroisses, des filtres hiérarchiques par région et district, un
tableau de bord adaptatif par rôle, un module d'import/export des données
institutionnelles, et un journal d'audit automatique et exhaustif de toutes
les opérations d'écriture. L'ensemble des six objectifs formulés en
Introduction Générale a été atteint. Par rapport à la plateforme
GeoEEC~v1.0, la solution apporte une valeur ajoutée précise et documentée
sur cinq axes~: sécurité de l'authentification par sessions \textit{HttpOnly},
RBAC aligné sur les quatre niveaux de la hiérarchie synodale réelle,
publication cartographique conforme aux standards OGC WMS et WFS, traçabilité
automatique par journal d'audit, et reproductibilité de l'environnement
d'exécution par conteneurisation complète.

\medskip
\noindent\textbf{Limites.}
Quatre limites structurelles tempèrent la portée de ces résultats. Premièrement,
480~paroisses sur 858 (56~\%) ne disposent d'aucune coordonnée GPS valide~:
la plateforme les enregistre mais ne peut pas les afficher sur la carte, ce
qui réduit d'autant sa capacité à soutenir des décisions de planification
territoriale fondées sur la couverture cartographique réelle. Deuxièmement,
le calcul d'itinéraire réel dépend de services de routing externes (OSRM,
Valhalla) non hébergés par l'institution~: en cas d'indisponibilité
prolongée de ces services publics, la plateforme retombe sur une distance
à vol d'oiseau (Haversine), moins précise pour les décisions opérationnelles.
Troisièmement, l'interface d'administration classique
n'est pas adaptée aux smartphones, qui constituent le principal outil
numérique des administrateurs paroissiaux en zone rurale~; la collecte
terrain des coordonnées GPS manquantes reste donc largement manuelle.
Quatrièmement, l'authentification ne propose pas de double facteur~: les
comptes d'administrateurs nationaux, qui disposent d'un accès complet à
toutes les données de l'institution, restent exposés aux risques de
compromission par mot de passe seul.

\medskip
\noindent\textbf{Perspectives.}
Ces limites définissent directement les axes de développement prioritaires.
La première perspective est le développement d'une application mobile
légère (React Native ou Progressive Web App) permettant aux administrateurs
paroissiaux de saisir ou de corriger les coordonnées GPS directement sur le
terrain~: cette seule fonctionnalité permettrait de résorber le déficit de
géolocalisation des 480~paroisses non couvertes et de faire passer le taux
de couverture cartographique de 44~\% à 100~\%. La deuxième perspective est
l'hébergement d'une instance propre de Valhalla ou d'OSRM sur
l'infrastructure de l'institution, afin de s'affranchir de la dépendance
aux services publics identifiée ci-dessus tout en conservant la même
famille d'algorithmes (section~\ref{sec:modelisation-math}). La troisième
perspective est l'activation de la double authentification par TOTP
(\textit{Time-based One-Time Password}) pour les comptes d'administrateurs
nationaux et régionaux, en s'appuyant sur la bibliothèque
\texttt{django-otp} déjà compatible avec l'infrastructure Django. Enfin, la
réintégration du module de questions-réponses en langage naturel développé
dans GeoEEC~v1.0 --- enrichi des nouvelles entités, des classes de
paroisses et des statistiques annuelles --- ouvrirait la plateforme à des
usages analytiques avancés permettant aux responsables synodaux
d'interroger les données institutionnelles sans compétences techniques.

\noindent\rule{\linewidth}{0.8pt}

\clearpage

% ============================================================
%  RÉFÉRENCES
% ============================================================
\chapter*{Références}
\addcontentsline{toc}{chapter}{Références}
\markboth{Références}{Références}

\noindent
Le tableau ci-dessous recense, par ordre de première apparition dans le
document, l'ensemble des sources citées dans le texte sous forme numérique
--- \textbf{[1]}, \textbf{[2]}, \textbf{[3]}, etc. Chaque source conserve le
même numéro à chacune de ses citations ultérieures dans le rapport.

\vspace{6pt}

\setlength{\tabcolsep}{6pt}
\renewcommand{\arraystretch}{1.3}
\begin{longtable}{%
  >{\bfseries\color{eccdarkgreen}\centering\arraybackslash}p{1.6cm}
  >{\raggedright\arraybackslash}p{11.4cm}}
\caption{Références bibliographiques du rapport}\label{tab:references} \\
\toprule
\rowcolor{eecgreen}
\multicolumn{1}{>{\bfseries\color{white}\centering\arraybackslash}p{1.6cm}}{Numéro} &
\multicolumn{1}{>{\bfseries\color{white}}p{11.4cm}}{Référence bibliographique} \\
\midrule
\endfirsthead
\caption[]{\textit{(suite)~Références bibliographiques du rapport}} \\
\toprule
\rowcolor{eecgreen}
\multicolumn{1}{>{\bfseries\color{white}\centering\arraybackslash}p{1.6cm}}{Numéro} &
\multicolumn{1}{>{\bfseries\color{white}}p{11.4cm}}{Référence bibliographique} \\
\midrule
\endhead
\midrule
\multicolumn{2}{r}{\small\textit{(suite page suivante)}} \\
\endfoot
\bottomrule
\endlastfoot
[1] & International Telecommunication Union, \textit{Measuring Digital Development: Facts and Figures 2023}, ITU, Genève, 2023. Disponible sur~: \url{https://www.itu.int/en/ITU-D/Statistics/Pages/facts/default.aspx} \\
[2] & GSMA Intelligence, \textit{The Mobile Economy Sub-Saharan Africa 2023}, GSM Association, Londres, 2023. Disponible sur~: \url{https://www.gsma.com/mobileeconomy/sub-saharan-africa/} \\
[3] & République du Cameroun, \textit{Stratégie Nationale de Développement 2020--2030~: Pour la transformation structurelle et le développement inclusif}, Ministère de l'Économie, de la Planification et de l'Aménagement du Territoire, Yaoundé, 2020. \\
[4] & P.~A. Longley, M.~F. Goodchild, D.~J. Maguire, D.~W. Rhind, \textit{Geographic Information Science and Systems}, 4\textsuperscript{e} édition, Wiley, Hoboken, NJ, 2015. \\
[5] & Open Geospatial Consortium, \textit{OpenGIS Web Map Service (WMS) Implementation Specification}, version~1.3.0, OGC~06-042, 2006. Disponible sur~: \url{https://www.ogc.org/standards/wms} \\
[6] & Open Geospatial Consortium, \textit{OpenGIS Web Feature Service (WFS) Implementation Standard}, version~2.0.0, OGC~09-025r2, 2010. Disponible sur~: \url{https://www.ogc.org/standards/wfs} \\
[7] & P.~Gifford, «~Christianity, Development and Modernity in Africa~», \textit{African Affairs}, vol.~114, n\textsuperscript{o}~454, p.~1--19, 2015. \\
[8] & Esri, \textit{ArcGIS Online: Cloud-Based Mapping and Analysis}, 2024. Disponible sur~: \url{https://www.esri.com/en-us/arcgis/products/arcgis-online/overview} \\
[9] & Google LLC, \textit{Google Maps Platform Documentation}, 2024. Disponible sur~: \url{https://developers.google.com/maps} \\
[10] & QGIS Development Team, \textit{QGIS Geographic Information System}, Open Source Geospatial Foundation Project, 2024. Disponible sur~: \url{https://www.qgis.org} \\
[11] & OWASP Foundation, \textit{OWASP Top Ten 2021: The Ten Most Critical Web Application Security Risks}, 2021. Disponible sur~: \url{https://owasp.org/www-project-top-ten/} \\
[12] & PostGIS Development Team, \textit{PostGIS 3.4 Documentation}, 2023. Disponible sur~: \url{https://postgis.net/documentation/} \\
[13] & R.~T. Fielding, \textit{Architectural Styles and the Design of Network-based Software Architectures}, thèse de doctorat, University of California, Irvine, 2000. \\
[14] & T.~Christie, \textit{Django REST Framework Documentation}, 2024. Disponible sur~: \url{https://www.django-rest-framework.org} \\
[15] & Docker, Inc., \textit{Docker Documentation}, 2024. Disponible sur~: \url{https://docs.docker.com} \\
[16] & E.~W. Dijkstra, «~A note on two problems in connexion with graphs~», \textit{Numerische Mathematik}, vol.~1, n\textsuperscript{o}~1, p.~269--271, 1959. \\
[17] & P.~E. Hart, N.~J. Nilsson, B.~Raphael, «~A Formal Basis for the Heuristic Determination of Minimum Cost Paths~», \textit{IEEE Transactions on Systems Science and Cybernetics}, vol.~4, n\textsuperscript{o}~2, p.~100--107, 1968. \\
[18] & Django Software Foundation, \textit{Django Documentation}, version~5.0, 2024. Disponible sur~: \url{https://docs.djangoproject.com/en/5.0/} \\
[19] & GeoServer Development Team, \textit{GeoServer 2.26 User Manual}, 2024. Disponible sur~: \url{https://docs.geoserver.org/stable/en/user/} \\
[20] & Vercel, \textit{Next.js Documentation}, version~16, 2024. Disponible sur~: \url{https://nextjs.org/docs} \\
[21] & V.~Agafonkin, \textit{Leaflet: An open-source JavaScript library for mobile-friendly interactive maps}, 2024. Disponible sur~: \url{https://leafletjs.com} \\
[22] & Redis Ltd., \textit{Redis Documentation}, 2024. Disponible sur~: \url{https://redis.io/docs} \\
[23] & OSRM Contributors, \textit{Project OSRM --- Open Source Routing Machine}, 2024. Disponible sur~: \url{https://project-osrm.org} \\
[24] & Valhalla Contributors, \textit{Valhalla Routing Engine Documentation}, FOSSGIS, 2024. Disponible sur~: \url{https://valhalla.github.io/valhalla/} \\
[25] & MapLibre Organization, \textit{MapLibre GL JS Documentation}, 2024. Disponible sur~: \url{https://maplibre.org/maplibre-gl-js/docs/} \\
\end{longtable}

\end{document}
